\label{C:section}

We use exact representation groups in C.1 and split quadratic-form groups in
C.3.  These are different conventions.  Cote 133, scans 7 and 28--29, distinguishes
$K_s$ from $K$; scan 6 defines $WG$ using isometry classes and orthogonal sums
only.  We inspected the official images of scans 6, 26, and 36 in this run.

\subsection{C.1: exact representations of symmetric groups}

For a commutative unital ring $k$, let $\mathcal R_k(G)$ be the exact category
of left $k[G]$-modules which are finitely generated projective over $k$; its
admissible sequences are short exact sequences of $k[G]$-modules.  Put
$R_k(G)=K_0(\mathcal R_k(G))$.  Thus every admissible sequence splits over $k$,
but it need not split $G$-equivariantly.  Induction and the external tensor
product give a graded ring
\[
 R_k(\mathfrak S_*)=\bigoplus_{n\geq0}R_k(\mathfrak S_n),\qquad
 [U]*[V]=[\operatorname{Ind}_{\mathfrak S_a\times\mathfrak S_b}^{\mathfrak S_{a+b}}
 (U\otimes_k V)].
\]
Let $\rho_n=[k]$ in degree $n$, with trivial action, and let
$\vartheta_k:\mathbb Z[h_1,h_2,\ldots]\longrightarrow R_k(\mathfrak S_*)$
send $h_n$ to $\rho_n$.

\begin{proposition}[Proved: the complete answer for $\mathbb Q$-algebras]
\label{C:Qalgebras}
For every commutative $\mathbb Q$-algebra $k$, without a noetherian or reduced
hypothesis, there is a natural graded-ring isomorphism
\[
 R_k(\mathfrak S_*)\simeq K_0(k)\otimes_{\mathbb Z}\mathbb Z[h_1,h_2,\ldots].
\]
Consequently $\vartheta_k$ is an isomorphism if and only if
$\mathbb Z\longrightarrow K_0(k)$, $1\mapsto[k]$, is an isomorphism.
In particular this holds for every nonzero local $\mathbb Q$-algebra.
\end{proposition}
\begin{proof}
The split Wedderburn decomposition of $\mathbb Q[\mathfrak S_n]$ gives
$k[\mathfrak S_n]\simeq\prod_{\lambda\vdash n}M_{d_\lambda}(k)$ after tensoring
with $k$.  The associated matrix idempotents decompose every object of
$\mathcal R_k(\mathfrak S_n)$ into
$\bigoplus_\lambda V_\lambda\otimes_{\mathbb Q}P_\lambda$, where each
$P_\lambda$ is a finite projective $k$-module. Conversely every such collection
gives an object. These exact equivalences identify the degree-$n$ group with
$\bigoplus_{\lambda\vdash n}K_0(k)$. The multiplicities for induction are the
integer multiplicities over $\mathbb Q$, since the decomposition over $\mathbb Q$
can be tensored with arbitrary $P_\lambda\otimes_kP_\mu$.
The rational-character induction ring is the integral symmetric-function ring,
with trivial representations corresponding to the complete symmetric functions;
equivalently, the Young permutation characters and the irreducible characters
are related by a unitriangular Kostka matrix. This proves the displayed
isomorphism, including multiplication. Its degree-zero part is exactly the
specified map $\mathbb Z\to K_0(k)$, proving necessity and sufficiency.
Finite projective modules over a local ring are free, proving the final claim.
\end{proof}


\begin{proposition}[Proved: every degree away from its finitely many bad primes]
\label{C:tame_degree}
Let $d\geq0$ and let $k$ be any commutative ring in which $d!$ is
invertible. For $0\leq n\leq d$,
\[
 R_k(\mathfrak S_n)\simeq\bigoplus_{\lambda\vdash n}K_0(k).
\]
These identifications respect induction whenever the total degree is at
most $d$. Equivalently, after discarding products in degrees above $d$,
the induction ring is $K_0(k)\otimes\Lambda_{\leq d}$ with its usual
symmetric-function structure. Thus for any local ring with residue
characteristic $p$, without reducedness or noetherianity, the polynomial
description always holds in every degree strictly smaller than $p$.
\end{proposition}
\begin{proof}
Put $A=\mathbb Z[1/d!]$. For $n\leq d$, let $S_A^\lambda$ be the free
integral cellular Specht lattices specialized to $A$. Their action gives
an $A$-algebra map
\[
 A[\mathfrak S_n]\longrightarrow
       \prod_{\lambda\vdash n}\operatorname{End}_A(S_A^\lambda).
\]
Both sides are finite free of the same rank: the equality
$n!=\sum_\lambda(\operatorname{rank}S_A^\lambda)^2$ is the count of
the cellular basis in Dipper--James--Mathas,
\cite[Theorem 3.26]{C-DipperJamesMathas1998}. Over every residue field
of $A$, averaging by $1/n!$ gives semisimplicity. All cell modules are
then absolutely irreducible and form the complete simple list, by
Graham--Lehrer \cite[Proposition 2.4(i),(v)]{C-GrahamLehrer2004}.
The displayed map is consequently an isomorphism on every residue
field. Its determinant is a unit of $A$, so it is an isomorphism over
$A$ and remains one after arbitrary base change to $k$.
Morita equivalence now gives the displayed representation groups.

The induction multiplicity modules over the PID $A$ are finite
projective, hence free; their ranks are the rational induction
multiplicities. Base change to $k$ gives the required products.
The unitriangular change between the Specht and Young permutation
bases, equivalently Burroughs's symmetric-function theorem
\cite[Theorem 1.2]{C-Burroughs1974}, identifies the latter generators
with the $\rho_i$. Finally every integer prime to $p$ is a unit in a
local ring of residue characteristic $p$, and $K_0(k)=\mathbb Z$.
This proves the last assertion, including all nilpotent and
nonnoetherian local bases in that degree range.
\end{proof}

\begin{proposition}[Proved: a universal obstruction in positive characteristic]
\label{C:positive_characteristic}
Let $k\ne0$ have positive characteristic $N$. Then $N\geq2$ and
$\vartheta_k$ is not injective. More explicitly, the nonzero Newton power-sum
polynomial $p_N(h_1,\ldots,h_N)$ maps to zero in $R_k(\mathfrak S_N)$.
This includes rings of characteristic $p^r$, not just fields of characteristic
$p$.
\end{proposition}
\begin{proof}
Let $V=k^N$ with its permutation action and let
$W=\ker(\epsilon:V\to k)$, where $\epsilon(x)=\sum_i x_i$. The vector
$w=(1,\ldots,1)$ belongs to $W$, since $N=0$ in $k$.
The vectors $e_i-e_N$ for $1\leq i<N$ are a basis of $W$, and in that basis
$w$ has all coordinates equal to $1$. Hence $kw\subset W$ is a trivial
rank-one subrepresentation with free quotient. Exterior powers of an admissible
sequence have their usual finite filtration by exterior powers of its ends;
all terms are finite projective over $k$. In the exact representation ring,
\[
 \lambda_{-1}(W)=\lambda_{-1}(k)\lambda_{-1}(W/kw)=0.
\]
Writing $e_j$ for the sign class in degree $j$, we have
$[\Lambda^jV]=e_j\rho_{N-j}$, by the signed-subset description of exterior
powers of a permutation module. From $0\to W\to V\to k\to0$ we get
\[
 \lambda_{-1}(W)
 =\sum_{j=0}^{N-1}(-1)^j(N-j)e_j\rho_{N-j}
 =p_N(\rho_1,\ldots,\rho_N).
\]
Here the relations expressing $e_j$ in the $\rho_i$ follow from the augmented
simplex complex on a $j$-element set: its equivariant Koszul differential is
contraction by the sum-of-coordinates map, and the complex is exact over
$\mathbb Z$, with free kernels. Thus these relations survive every base
change. They are precisely $e(t)h(-t)=1$, and the last equality is Newton's
identity. The coefficient of $h_N$ in $p_N$ is $N$, so the polynomial is
nonzero over $\mathbb Z$. The case $N=2$ is the checked relation
$\rho_1^2=2\rho_2$ from the supplied Pro notes F11.
\end{proof}

\begin{corollary}[Proved: injectivity is completely classified for all rings]
\label{C:allrings_injectivity}
For a commutative ring $k$, the map $\vartheta_k$ is injective if and only
if $k$ is nonzero of characteristic zero. In that case it is a split
injection of graded rings after choosing a characteristic-zero residue
field of $k\otimes\mathbb Q$. Thus the remaining all-rings problem is
entirely a generation problem, not a question about hidden relations
among the $\rho_i$.
\end{corollary}
\begin{proof}
The zero ring has zero target, and positive characteristic was excluded
in Proposition~\ref{C:positive_characteristic}. If $\operatorname{char}k=0$,
then $B=k\otimes\mathbb Q$ is nonzero, since $1\otimes1=0$ would imply
that some nonzero integer annihilates $1\in k$. Choose a maximal ideal
of $B$ and let $F$ be its residue field. Scalar extension preserves all
admissible sequences, which are split on the underlying $k$-modules,
and gives $R_k(\mathfrak S_*)\to R_F(\mathfrak S_*)$. Its composite
with $\vartheta_k$ is the isomorphism $\vartheta_F$ from
Proposition~\ref{C:Qalgebras}. Composing with $\vartheta_F^{-1}$ gives
the required ring retraction.
\end{proof}

\begin{proposition}[Refuted: locality, even with torsion-free characteristic zero]
\label{C:local_charzero_counterexample}
The ring
\[
 k=\mathbb Z_{(2)}[u]/(u^2-1)
\]
is nonzero, local, of characteristic zero, and torsion-free over $\mathbb Z$;
$K_0(k)=\mathbb Z$. Nevertheless $\vartheta_k$ is not surjective in degree two.
\end{proposition}
\begin{proof}
The ring is free with basis $1,u$ over $\mathbb Z_{(2)}$, and is finite integral
over that local ring. Every maximal ideal therefore lies over $(2)$.
Its reduction is $\mathbb F_2[u]/((u-1)^2)$, with exactly one maximal ideal.
Thus $k$ is local, and finite projective $k$-modules are free. The basis also
proves characteristic zero and torsion-freeness.
Let $L$ be the free rank-one $k$-module on which the nonidentity element of
$\mathfrak S_2$ acts as multiplication by $u$.
Taking determinants of this element gives a group homomorphism
\[
 \det_s:R_k(\mathfrak S_2)\longrightarrow k^\times,
\]
because determinant is multiplicative on every admissible exact sequence.
It takes $\rho_2$ to $1$ and $\rho_1^2$ (the regular representation) to $-1$.
Consequently every element in the degree-two image of $\vartheta_k$ has
determinant $1$ or $-1$. But $\det_s(L)=u$, and the basis $1,u$ proves
$u\ne1,-1$. This excludes $[L]$ from that image. The determinant obstruction
can already be distinguished in the finite quotient
$(\mathbb Z/4)[u]/(u^2-1)$.
\end{proof}

\begin{proposition}[Proved: an obstruction invisible on the rational fibre]
\label{C:torsion_generic_kernel}
For every prime $p$, let
\[
 k=\mathbb Z_{(p)}[\epsilon]/(p\epsilon,\epsilon^p).
\]
This is a local ring of characteristic zero, with $K_0(k)=\mathbb Z$ and
$k\otimes\mathbb Q=\mathbb Q$. In fact $K_0(k[1/N])=\mathbb Z$ for
every nonzero integer $N$. Nevertheless $\vartheta_k$ is not surjective
in degree $p$. The comparison
$R_k(\mathfrak S_p)\to R_{\mathbb Q}(\mathfrak S_p)$ has a nonzero
kernel, so this obstruction is distinct from the generic Picard class in
Proposition~\ref{C:normal_domain_counterexample}.
\end{proposition}
\begin{proof}
As an abelian group $k$ is
$\mathbb Z_{(p)}\oplus\bigoplus_{i=1}^{p-1}\mathbb F_p\epsilon^i$.
In particular $\epsilon^{p-1}\ne0$ and is not a constant. The nilpotent
ideal $(\epsilon)$ has local quotient $\mathbb Z_{(p)}$, so $k$ is local.
If $p\nmid N$, localization at $N$ changes nothing; if $p\mid N$, it
gives $\mathbb Q$. The assertions about characteristic and $K_0$ follow.

The unit $u=1+\epsilon$ satisfies $u^p=1$: every intermediate binomial
coefficient is divisible by $p$, and $\epsilon^p=0$. Let $C_p\subset
\mathfrak S_p$ be generated by a $p$-cycle $g$, and let $L$ be the free
rank-one $k$-representation with $g$ acting by $u$. Put
$M=\operatorname{Ind}_{C_p}^{\mathfrak S_p}L$ and
$M_0=\operatorname{Ind}_{C_p}^{\mathfrak S_p}k$. The cosets fixed by $g$
in this induction are precisely the cosets of $C_p$ in its normalizer;
that quotient acts on $C_p$ through all $a\in\mathbb F_p^\times$.
The induction trace formula therefore gives
\[
 \operatorname{tr}(g\mid M)
   =\sum_{a=1}^{p-1}(1+\epsilon)^a
   =p-1+\epsilon^{p-1}.
\]
For the last equality, the coefficient of $\epsilon^i$ in the sum
including $a=0$ is $\binom p{i+1}$, by the elementary binomial-sum
identity. All coefficients with $1\leq i<p-1$ vanish in $k$.

Trace at $g$ is additive on admissible exact sequences of finite
projective $k$-modules. Every degree-$p$ monomial in the $\rho_i$ is a
Young permutation representation, whose trace is an integer. Thus the
displayed trace excludes $[M]$ from the image of $\vartheta_k$.
Moreover $M$ and $M_0$ become isomorphic over $\mathbb Q$, while their
trace difference over $k$ is the nonzero $\epsilon^{p-1}$. Their
difference gives the asserted nonzero generic-comparison kernel.

For $p=3$ the inducing representation has dimension two, with matrices
\[
 g\longmapsto\begin{pmatrix}u&0\\0&u^2\end{pmatrix},\qquad
 s\longmapsto\begin{pmatrix}0&1\\1&0\end{pmatrix}
\]
for a three-cycle and a transposition. The attached program
\path{code/C_torsion_trace_certificate.py} verifies the complete
$S_3$ multiplication table and the nonconstant trace in the 81-element
quotient $(\mathbb Z/9)[\epsilon]/(3\epsilon,\epsilon^3)$; it also checks
the ring axioms on every triple of that quotient.
\end{proof}

\begin{proposition}[Refuted: even normal local domains of characteristic zero]
\label{C:normal_domain_counterexample}
Let
\[
 R=\left(\mathbb Z_{(2)}[a,b,c]/(bc-a(2-a))\right)_{(2,a,b,c)}.
\]
Then $R$ is a normal noetherian local domain of characteristic zero,
$K_0(R)=\mathbb Z$, and $R\otimes\mathbb Q=R[1/2]$ is a connected regular
ring. Nevertheless $\vartheta_R$ is not surjective in degree two.
\end{proposition}
\begin{proof}
The polynomial $bc-a(2-a)$ is primitive and irreducible as a degree-one
polynomial in $c$ over the UFD $\mathbb Z_{(2)}[a,b]$: its two coefficients
$b$ and $-a(2-a)$ are relatively prime. Thus $R$ is a domain of
characteristic zero. It is a hypersurface in a regular local ring, hence
Cohen--Macaulay and satisfies $(S_2)$. After inverting two its defining
hypersurface is smooth: the three partial derivatives are $2a-2,c,b$,
and their simultaneous vanishing would give $a=1,b=c=0$, inconsistent
with the equation. The ideal $(2)$ is prime, since its quotient is the
local domain $\mathbb F_2[a,b,c]/(a^2+bc)$. The quadratic polynomial has
no linear factor; this is also checked by the attached exact program.
Consequently the only height-one prime containing two is $(2)$ itself,
and $R_{(2)}$ is regular of dimension one with parameter two. Thus $R$
satisfies $(R_1)$ and is normal by Serre's criterion
\cite[Tag 031S]{C-StacksNormal}. The stated connectedness and regularity of
$B=R[1/2]$ follow at once, and locality gives $K_0(R)=\mathbb Z$.

Put $I=(a,b)$. This is a height-one prime: the quotient is
$\mathbb Z_{(2)}[c]_{(2,c)}$, of dimension two, while $\dim R=3$.
It is not principal. Indeed the hypersurface equation lies in the square
of the ambient maximal ideal, so the images of $2,a,b,c$ are linearly
independent in $\mathfrak m/\mathfrak m^2$. In particular $a,b$ remain
independent in $I/\mathfrak m I$, and Nakayama's lemma says that $I$ needs
two generators. In a normal domain a height-one prime is the intersection
of its localizations at height-one primes, inside the fraction field:
at $I$ this follows by contraction of the prime, and at every other
height-one prime its localization is the whole ring. The intersection
statement for $R$ is \cite[Tag 031T]{C-StacksNormal}.

We claim that $IB$ is still not principal. If $IB=xB$, let $v_{(2)}(x)=m$
and replace $x$ by $y=2^{-m}x$. For every height-one prime other than $(2)$,
localization has already inverted two and gives $I R_{\mathfrak p}=y
R_{\mathfrak p}$. At $(2)$, $I R_{(2)}=R_{(2)}$ because $a\notin(2)$,
and the choice of $y$ gives the same equality. Intersecting these equalities
at all height-one primes forces $I=yR$, a contradiction. This argument
requires no computation of the entire divisor class group.

Now let the transposition in $\mathfrak S_2$ act on $R^2$ by
\[
 T=\begin{pmatrix}a-1&b\\c&1-a\end{pmatrix}.
\]
Its square is the identity, since $T^2-1=(bc-a(2-a))1_2$.
Over $B$ the idempotent $e=(T+1)/2$ has rank one. Projection to the first
coordinate identifies $\operatorname{im}(e)$ with $IB$: its image is
$(a,b)B$, and its kernel is zero because it is a map from a rank-one
torsion-free module that is nonzero generically. Thus $IB$ is invertible
but nontrivial in $\operatorname{Pic}(B)$.

The exact functor taking the positive eigenspace after extension to $B$
induces $R_R(\mathfrak S_2)\to K_0(B)$; composing with determinant gives
an additive map to $\operatorname{Pic}(B)$. Both $\rho_2$ and $\rho_1^2$
have a free positive eigenspace, so their images are zero. The representation
$T$ has image $[IB]\ne0$. Hence its class is outside the degree-two image
of $\vartheta_R$, as claimed. The symbolic matrix identity and the
mod-two irreducibility check appear in
\texttt{code/C\_involution\_certificate.py} and its JSON output.
\end{proof}

\begin{proposition}[Known for $\mathbb Z$; proved reduction for mixed-characteristic DVRs]
\label{C:DVR}
The map $\vartheta_{\mathbb Z}$ is an isomorphism. It is also an isomorphism
when $k$ is any discrete valuation ring of mixed characteristic, in particular
when $k=\mathbb Z_p$ for any prime $p$.
\end{proposition}
\begin{proof}
The integral assertion is Burroughs, Proposition 1.1 and Theorem 1.2
\cite{C-Burroughs1974}. We do not use the unrestricted Dedekind-ring paraphrase
in that proposition's proof: for a general Dedekind ring even the trivial
group has $K_0(k)=\mathbb Z\oplus\operatorname{Pic}(k)$.

For the DVR assertion let $K$ be the fraction field, $\kappa$ the residue
field, and $\varpi$ a uniformizer. Every finitely generated $k[\mathfrak S_n]$
module admits a surjection from a finite free group-ring module whose kernel
is torsion-free over $k$ and hence free over $k$. For clarity, this identifies $R_k(\mathfrak S_n)$ with
$G_0(k[\mathfrak S_n])$ without an unproved resolution comparison:
for a presentation $0\to L\to A^r\to M\to0$, where $A=k[\mathfrak S_n]$,
assign $[M]\mapsto[A^r]-[L]$. Schanuel's lemma makes this independent of
the presentation, and the horseshoe construction makes it additive on
exact sequences. It is visibly inverse to the natural map on lattice
classes.
The localization sequence is
\[
 G_0(\kappa[\mathfrak S_n])\longrightarrow
 G_0(k[\mathfrak S_n])\longrightarrow G_0(K[\mathfrak S_n])\longrightarrow0.
\]
Here surjectivity follows by taking the $A$-span of a finite $K$-basis
inside any $K[\mathfrak S_n]$-module. Two such lattices are commensurable,
so their difference comes from finite $k$-torsion modules. Exact sequences
over $K$ lift by intersecting a lattice in the middle term with the submodule
and taking its image in the quotient. Finally, filtering a finite torsion
module by powers of $\varpi$ reduces it to $\kappa[\mathfrak S_n]$-modules.
These observations prove the displayed exactness at $G_0$ directly.
The left group is generated by reductions of integral Specht modules.
Here is the precise input, valid also at $p=2$: specialize
Dipper--James--Mathas \cite[Theorem 3.26, Definition 3.28, Theorem 3.30;
preprint pp.12--13]{C-DipperJamesMathas1998} to $r=1,q=1$.
The integral cellular basis gives free integral Specht lattices whose
specializations are the cell modules over every residue field; their
nonzero simple heads give all simples. The decomposition matrix has a
unitriangular submatrix, by Graham--Lehrer
\cite[Proposition 2.4(i),(ii); author manuscript p.3]{C-GrahamLehrer2004}.
Consequently every simple class is an integral combination of reductions
of these integral lattices. The right-module convention in the source
is changed to left modules by $g\mapsto g^{-1}$; using dual Specht
lattices instead has the same generation consequence. For each such
Specht lattice $S_k^\lambda$,
\[
 0\longrightarrow S_k^\lambda\xrightarrow{\varpi}S_k^\lambda
 \longrightarrow S_\kappa^\lambda\longrightarrow0
\]
shows that its reduction has zero image in the middle group. Thus the first
map is zero and passage to $K$ is an isomorphism. It commutes with induction
and sends every $\rho_n$ to $\rho_n$. Since $K$ has characteristic zero,
Proposition~\ref{C:Qalgebras} concludes the proof.
The Specht-generation and localization inputs are recorded separately in the
reference/check register; this argument requires both, not just a comparison
of ranks.
\end{proof}

\begin{lemma}[Proved: integral Specht generators over a noetherian base]
\label{C:Specht_G0}
Let $k$ be any commutative noetherian ring and $n\geq0$.
The homomorphism
\[
 \bigoplus_{\lambda\vdash n}G_0(k)\longrightarrow G_0(k[\mathfrak S_n]),
 \qquad ([M_\lambda])_\lambda\longmapsto
       \sum_\lambda[S^\lambda_{\mathbb Z}\otimes_{\mathbb Z}M_\lambda]
\]
is surjective. Either the usual integral Specht lattices or their duals
may be used consistently.
\end{lemma}
\begin{proof}
We use the fieldwise input that the classes of Specht modules span
$G_0(F[\mathfrak S_n])$ for every field $F$, including characteristic two.
It follows from the unitriangular submatrix of the cellular decomposition
matrix, with diagonal entries one, as checked in the proof of
Proposition~\ref{C:DVR}.

Put $A=k[\mathfrak S_n]$; it is noetherian and finite free over its central
subring $k$. We prove the assertion by noetherian induction on the closed
support of a finite $A$-module $M$. Such a module admits a finite filtration
by $A$-submodules with each quotient annihilated by a prime of $k$: in a
nonzero quotient, choose a $k$-associated element $x$ whose annihilator is
a prime $\mathfrak p$ and take the $A$-submodule $Ax$. Centrality implies
that $\mathfrak p$ annihilates $Ax$, and noetherianity prevents an infinite
strictly increasing sequence of these submodules. Thus it suffices to treat
$M$ over the domain $k/\mathfrak p$ with support $V(\mathfrak p)$.

Let $F=\operatorname{Frac}(k/\mathfrak p)$. The fieldwise input writes
$[M_F]=\sum_\lambda a_\lambda[S^\lambda_F]$. This equality involves only
finitely many short exact sequences and isomorphisms of finite-dimensional
modules. Their matrices and relations can be spread out after inverting
one $f\notin\mathfrak p$; kernels and cokernels vanish after a further such
localization because they are finite. Therefore
\[
 [M]-\sum_\lambda a_\lambda[S^\lambda_{\mathbb Z}\otimes(k/\mathfrak p)]
\]
is a sum of classes of modules supported on
$V(\mathfrak p)\cap V(f)$, a proper closed subset of $V(\mathfrak p)$.
Here the elementary $G_0$ localization assertion can be verified directly:
finite modules over a central localization have finite lattices, two lattices
are commensurable up to a power of the inverted element, and exact sequences
are lifted by intersecting a lattice with the submodule and taking its
image in the quotient. Thus the kernel is generated by modules killed by
powers of $f$. The induction hypothesis treats these smaller supports.
The displayed Specht terms already belong to the asserted image, with
coefficient $[k/\mathfrak p]\in G_0(k)$. This completes the induction.
\end{proof}

\begin{theorem}[Proved: the complete answer for regular noetherian rings]
\label{C:regular_noetherian}
Let $k\ne0$ be a regular noetherian commutative ring. Then
$\vartheta_k$ is an isomorphism if and only if $k$ has characteristic zero
and the unit map $\mathbb Z\to K_0(k)$ is an isomorphism. No bound on the
Krull dimension is required. In particular it is an isomorphism for every
regular noetherian local ring of characteristic zero.
\end{theorem}
\begin{proof}
Necessity of the $K_0$ condition follows in degree zero, and positive
characteristic is excluded by Proposition~\ref{C:positive_characteristic}.
Assume the two conditions. Every finite $k$-module has a finite resolution
by finite projective $k$-modules. To recall why an unbounded global dimension
causes no gap here, fix a free resolution of that one module. At each prime,
regularity makes some syzygy free; the loci where a given finite syzygy is
locally free are open and increase with the syzygy index. They cover the
quasi-compact spectrum, so a single finite index works everywhere.

Apply the same construction using finite free $k[\mathfrak S_n]$-modules:
a sufficiently high syzygy is finite projective over $k$. Alternating sums
of these resolution terms give an inverse
$G_0(k[\mathfrak S_n])\to R_k(\mathfrak S_n)$ to the natural map.
Independence follows by iterated Schanuel comparison, and additivity by
horseshoe resolutions; the terminal syzygies are projective over $k$, so
their comparison sequences are admissible in $\mathcal R_k(\mathfrak S_n)$.
This is the same explicit Euler-characteristic argument as for the DVR,
with a finite number of additional syzygies.

By Lemma~\ref{C:Specht_G0}, the resulting group is generated by the classes
$S^\lambda_{\mathbb Z}\otimes M$, for $M$ finite over $k$. Resolving $M$
over $k$ expresses each such class by Specht lattices tensored with finite
projective $k$-modules. The hypothesis $K_0(k)=\mathbb Z$ then expresses
all of them as integer combinations of $[S^\lambda_k]$.
Burroughs's integral theorem expresses $[S^\lambda_{\mathbb Z}]$ as a
polynomial in the $\rho_i$. Its exact sequences are split over $\mathbb Z$,
so arbitrary base change to $k$ preserves them. Thus $\vartheta_k$ is
surjective in every degree.

Since $k$ has characteristic zero, $k\otimes\mathbb Q\ne0$ and has a
residue field $F$ of characteristic zero. Base change $k\to F$ gives a
map $R_k(\mathfrak S_*)\to R_F(\mathfrak S_*)$ carrying the $\rho_i$ to
the corresponding generators. The composite with $\vartheta_k$ is the
isomorphism of Proposition~\ref{C:Qalgebras} for the field $F$, so
$\vartheta_k$ is injective. This completes the proof.
\end{proof}

\begin{theorem}[Proved: regularity is needed only along the modular locus]
\label{C:regular_modular_locus}
Let $k\ne0$ be noetherian of characteristic zero, with
$\mathbb Z\xrightarrow{\sim}K_0(k)$. Fix $n\geq0$. If $k_{\mathfrak p}$
is regular for every prime $\mathfrak p$ containing $n!$, then the
degree-$n$ map $\vartheta_k$ is an isomorphism.
Consequently, if $k_{\mathfrak p}$ is regular at every point of positive
residue characteristic, the entire map $\vartheta_k$ is an isomorphism.
Singular points of residue characteristic zero are allowed.
\end{theorem}
\begin{proof}
Put $t=n!$, $A=k[\mathfrak S_n]$, and $k_t=k[1/t]$.
We use the exact category $\mathcal C$ of finite $A$-modules of finite
projective dimension over $k$. Finite free $A$-resolutions, terminated
at a $k$-projective syzygy, identify $K_0(\mathcal C)$ with
$R_k(\mathfrak S_n)$ by the Euler-characteristic argument already
spelled out in Theorem~\ref{C:regular_noetherian}; regularity of all of
$k$ is unnecessary for modules in this particular category.

We first record the consequence of the stated regularity hypothesis.
Every finite $k$-module supported on $V(t)$ has finite projective
dimension: at each point of its support the local ring is regular,
and elsewhere the module is zero. The increasing open locally-free
loci of its finite syzygies and quasi-compactness give a uniform finite
resolution. The same argument shows that a finite $k$-module which is
projective after inverting $t$ has finite projective dimension over $k$.
Any finite projective $k_t$-module has a finite $k$-module lattice, so
its class lifts by this Euler resolution. Therefore
$K_0(k)\to K_0(k_t)$ is surjective. A characteristic-zero residue field
of $k\otimes\mathbb Q$ supplies a rank retraction; hence this map is an
isomorphism and both groups are $\mathbb Z$.

Consider
\[
 R_k(\mathfrak S_n)\longrightarrow R_{k_t}(\mathfrak S_n).
\]
Over $k_t$ every representation projective as a $k_t$-module is
projective as an $A_t$-module: average a $k_t$-linear splitting over
the group using $1/n!$. Choose a finite $A$-module lattice in any such
representation. Its underlying $k$-module has finite projective dimension
by the preceding paragraph, so its Euler class supplies a lift. Thus
the displayed map is surjective.

Its kernel is generated by the Euler classes of finite $A$-modules
supported on $V(t)$. Here is the comparison argument. Represent a kernel
class by $[M]-[N]$ with $M,N$ projective over $k$. Equality after
localization is stable isomorphism of finite projective $A_t$-modules,
so there is a finite projective $A_t$-module $Q$ with
$M_t\oplus Q\simeq N_t\oplus Q$. Lift $Q$ to a finite $A$-module
$L$ of finite $k$-projective dimension. Finite presentation permits
clearing denominators in the stabilized isomorphism to get
$f:M\oplus L\to N\oplus L$ that is an isomorphism after inverting
$t$. Its kernel and cokernel are finite, supported on $V(t)$, and
therefore have finite projective dimension over $k$. The image of $f$
also has finite projective dimension, by the two short exact sequences.
Their relations in $K_0(\mathcal C)$ give
$[M]-[N]=[\ker f]-[\operatorname{coker}f]$, as required.

Apply the support-induction proof of Lemma~\ref{C:Specht_G0} within
the abelian category of modules supported on $V(t)$. Every prime and
every smaller closed support used there stays inside $V(t)$. It
expresses their $G_0$ classes as sums of
$[S^\lambda_{\mathbb Z}\otimes T]$, where $T$ is a finite $k$-module
supported on $V(t)$. All these modules have finite $k$-projective
dimension, so the relations are valid in $K_0(\mathcal C)$ as well.
The Euler class of $T$ is zero in $K_0(k)$: it vanishes after
localization and $K_0(k)\to K_0(k_t)$ is an isomorphism.
Tensoring a finite projective resolution of $T$ with the free Specht
lattice therefore shows that every displayed class is zero in
$R_k(\mathfrak S_n)$. This proves injectivity of the comparison map.

Finally Proposition~\ref{C:tame_degree} identifies its target with
the degree-$n$ symmetric-function group, since $K_0(k_t)=\mathbb Z$.
The comparison preserves the $\rho_i$, proving the assertion in that
degree. Every prime containing some $n!$ has positive residue
characteristic, so the final hypothesis supplies this argument for
every $n$.
\end{proof}

\begin{proposition}[Proved: passage to filtered unions and limits of rings]
\label{C:filtered_colimits}
For a filtered system of commutative rings $k_i$ with colimit $k$, and
any finite group $G$, the natural map
\[
 \mathop{\rm colim}_i R_{k_i}(G)\longrightarrow R_k(G)
\]
is an isomorphism. For symmetric groups these isomorphisms respect induction
products. Consequently a filtered colimit of rings satisfying C.1 also
satisfies C.1. This gives, for example, nonnoetherian mixed-characteristic
local examples
$\mathbb Z_p[t_1,t_2,\ldots]_{(p,t_1,t_2,\ldots)}$.
\end{proposition}
\begin{proof}
A finite projective module is represented by a finite idempotent matrix.
An action of a finite group, an equivariant homomorphism, or an equivariant
isomorphism is specified by finitely many additional matrices and polynomial
identities. Each therefore descends to some $k_i$, and two such data become
equal at some later stage if they become equal over $k$.
An admissible short exact sequence is split on the underlying projective
modules; its equivariance, exactness, and a chosen $k$-linear splitting are
again finitely many matrix identities and descend as well. The presentation
of $K_0$ by objects and these finite relations therefore commutes with the
filtered colimit. The explicit finite matrices for induction commute with
base change, proving the ring assertion. For the example, take the filtered
system $\mathbb Z_p[t_1,\ldots,t_r]_{(p,t_1,\ldots,t_r)}$, whose stages
satisfy Theorem~\ref{C:regular_noetherian}.
\end{proof}

\begin{proposition}[Proved reduction: the two remaining obstruction groups]
\label{C:allrings_obstructions}
Let $k\ne0$ have characteristic zero, put $B=k\otimes\mathbb Q$, and
choose a residue field $F$ of $B$. Write $\Lambda=\mathbb Z[h_1,h_2,\ldots]$,
$J=\ker(K_0(B)\to K_0(F)=\mathbb Z)$, and
\[
 \beta:R_k(\mathfrak S_*)\longrightarrow
       R_B(\mathfrak S_*)=K_0(B)\otimes\Lambda.
\]
The retraction of Corollary~\ref{C:allrings_injectivity} has kernel $I$,
and there is an exact sequence of graded abelian groups
\[
 0\longrightarrow\ker\beta\longrightarrow I
   \xrightarrow{\beta}
      \operatorname{im}\beta\cap(J\otimes\Lambda)
   \longrightarrow0.                                  \tag{C.1.1}
\]
The induction ring is additively $\Lambda\oplus I$, with $I$ a graded
ideal. Hence $\vartheta_k$ is an isomorphism precisely when both end
obstruction groups in (C.1.1) vanish. These are independent obstructions:
Proposition~\ref{C:torsion_generic_kernel} exhibits a nonzero left group
with $J=0$, while Proposition~\ref{C:normal_domain_counterexample}
exhibits a nonzero right group over a torsion-free normal local domain.
\end{proposition}
\begin{proof}
The unit map $\mathbb Z\to K_0(B)$ splits the chosen rank map, so
$K_0(B)=\mathbb Z\oplus J$ as abelian groups. Tensoring with the free
abelian group $\Lambda$ identifies the kernel of the map to
$R_F(\mathfrak S_*)$ with $J\otimes\Lambda$. Taking its inverse image
under $\beta$ gives $I$ and proves exactness. The retraction and its
section are ring maps, so $I$ is an ideal and its additive complement
is $\vartheta_k(\Lambda)$. The last assertions follow from the indicated
counterexamples: the trace class dies on the rational fibre, whereas
the nontrivial positive eigenline survives in its $K_0$ but has rank
zero after subtracting a free line.
\end{proof}

\begin{lemma}[Proved: a Bockstein parity invariant of free involutions]
\label{C:Bockstein_parity}
Let $\mathcal E$ be the exact category of finite free
$\mathbb Z/4$-modules with an involution $T$. Set
$D=T-1$, $\overline M=M/2M$, and
\[
 H(M)=\ker(\overline D:\overline M\to\overline M)\big/
                         \operatorname{im}\overline D.
\]
There is a natural endomorphism $N_M$ of this $\mathbb F_2$-space:
if $\bar v$ is a cycle, choose a lift $v\in M$ and write $Dv=2w$;
then $N_M[\bar v]=[\bar w]$.
Let $f\in\mathbb F_2[t]$ be any monic irreducible polynomial satisfying
$f(t+1)=f(t)$, and let $H(M)_f$ denote the $f$-primary part
for $N_M$. Then
\[
 \chi_f(M)=\operatorname{length}_{\mathbb F_2[t]}H(M)_f\pmod2
       \quad\in\mathbb Z/2                              \tag{C.1.2}
\]
is additive on every admissible short exact sequence. It therefore
defines a homomorphism $R_{\mathbb Z/4}(C_2)\to\mathbb Z/2$.
\end{lemma}
\begin{proof}
The identity $D^2=-2D$ implies $\overline D^2=0$.
Because $M$ is free over $\mathbb Z/4$, division by $2$ determines
$\bar w$ uniquely. Also
$D w=D^2v/2=-Dv$ modulo $2M$, so $\bar w$ is a cycle.
Replacing $v$ by $v+2u$ changes $\bar w$ by the boundary
$\overline D\bar u$. If $\bar v=\overline D\bar z$ is a boundary,
choose $v=Dz$; then $Dv/2=-Dz$ modulo $2M$, also a boundary.
This proves well-definedness, linearity and naturality of $N_M$.

Consider an admissible sequence $0\to A\to E\to B\to0$ in
$\mathcal E$. An underlying $\mathbb Z/4$-linear splitting gives
\[
 T_E=\begin{pmatrix}T_A&X\\0&T_B\end{pmatrix},\qquad
 D_AX+XD_B+2X=0.                                      \tag{C.1.3}
\]
The short exact sequence of one-periodic complexes with differentials
$\overline D$ has the exact cyclic homology sequence
\[
 H(A)\longrightarrow H(E)\longrightarrow H(B)
            \xrightarrow{\partial}H(A).
                                                               \tag{C.1.4}
\]
Here $\partial[\bar b]=[\bar X\bar b]$, and exactness follows by
the ordinary cycle-and-boundary lifting argument. The first two
maps commute with $N$. For the last map write $D_Bb=2b'$ in (C.1.3);
division by $2$ gives
\[
 N_A\partial=\partial(N_B+1).                         \tag{C.1.5}
\]
In particular the connecting map need not commute with $N$.
This twist is essential to the invariant.

Since $f(t+1)=f(t)$, every map in (C.1.4) preserves the $f$-primary
parts. The restricted cyclic sequence is still exact: decompose each
space into its primary parts, observing that $\partial$ permutes
the parts by $g(t)\mapsto g(t+1)$; only the $f$-part can map to
the $f$-part. Write $a,e,b$ for the lengths of its three spaces,
and $r,s,t$ for the lengths of the three successive images.
Exactness gives $a=t+r$, $e=r+s$, $b=s+t$.
These lengths are legitimate also at the twisted connecting map:
twisting $t$ to $t+1$ preserves $f$-primary length. Consequently
$a+e+b=2(r+s+t)$, so $e\equiv a+b\pmod2$.
This proves the asserted additivity.
\end{proof}

\begin{proposition}[Proved: flat nilpotent mixed-characteristic bases can have a generic kernel]
\label{C:flat_nilpotent_kernel}
Let
\[
 k=\mathbb Z_{(2)}[\epsilon]/(\epsilon^2),\qquad
 T=\begin{pmatrix}1&\epsilon\\\epsilon&-1\end{pmatrix},
 \qquad M=(k^2,T).
\]
Then $k$ is local, flat over $\mathbb Z_{(2)}$, and has characteristic
zero. Every integer localization $k[1/a]$ with $a\ne0$ has $K_0=\mathbb Z$,
including $k\otimes\mathbb Q$.
Nevertheless
\[
 \eta=[M]-[\mathbf1]-[\mathrm{sgn}]
       \in R_k(\mathfrak S_2)                              \tag{C.1.6}
\]
is a nonzero element of order exactly two, killed by extension to
$k\otimes\mathbb Q$. In particular $\vartheta_k$ is not surjective.
The same assertion holds with $\mathbb Z_2$ in place of
$\mathbb Z_{(2)}$.
\end{proposition}
\begin{proof}
The ring is free of rank two over $\mathbb Z_{(2)}$ and local with
maximal ideal $(2,\epsilon)$. An integer localization is either this
ring or the local ring $\mathbb Q[\epsilon]/\epsilon^2$, so its
finite projectives are free and $K_0=\mathbb Z$.
The identity $T^2=1$ follows from $\epsilon^2=0$.
There is a ring homomorphism
\[
 k\longrightarrow\mathbb Z/4,\qquad\epsilon\longmapsto2.
\]
Scalar extension is exact on the representation category because
every admissible sequence splits on the underlying projective
modules. Compose it with Lemma~\ref{C:Bockstein_parity} for
$f=t^2+t+1$.
For the resulting involution on $(\mathbb Z/4)^2$ we obtain
\[
 \overline D=0,\qquad
 N_M=\begin{pmatrix}0&1\\1&1\end{pmatrix},
       \qquad\det(t-N_M)=t^2+t+1.
\]
Thus $H(M)$ is the simple $\mathbb F_2[t]/(f)$-module and
$\chi_f(M)=1$. For the trivial and sign lines, $N$ is respectively
$0$ and $1$, and $\chi_f=0$ in both cases. This proves $\eta\ne0$.
The degree-two image of $\vartheta_k$ is generated by the trivial
line and the regular permutation module. The exact sequence
$0\to\mathbf1\to k[C_2]\to\mathrm{sgn}\to0$ identifies this
with the span of the two lines, so the same invariant proves
failure of surjectivity.

After inverting $2$, the idempotents $(1\pm T)/2$ split $M$ into
rank-one projectives. The localized ring is local, so these are
the free trivial and sign lines, and $\eta$ becomes zero.
To compute its order, the standard evaluation sequence is
\[
 0\longrightarrow M\otimes\mathrm{sgn}
 \longrightarrow k[C_2]\otimes_k k^2
 \longrightarrow M\longrightarrow0.                    \tag{C.1.7}
\]
In the middle term $C_2$ acts on its first factor. Writing its
underlying module as $k^2\oplus k^2$, the last map is
$(u,v)\mapsto u+Tv$ and the first is $v\mapsto(-Tv,v)$;
the induced action on the kernel parameter is $-T$.
The matrix
$J=\left(\begin{smallmatrix}0&1\\-1&0\end{smallmatrix}\right)$
satisfies $JT=(-T)J$ and is invertible. Hence
$M\otimes\mathrm{sgn}\simeq M$ and (C.1.7) gives
$2[M]=2[k[C_2]]=2[\mathbf1]+2[\mathrm{sgn}]$.
Together with the nonzero invariant this proves that $\eta$ has
order exactly two. The proof for $\mathbb Z_2$ is identical, using
its quotient map to $\mathbb Z/4$.
\end{proof}

\begin{corollary}[Proved: infinitely many independent flat nilpotent obstructions]
\label{C:flat_nilpotent_infinite_kernel}
For either ring $k$ in Proposition~\ref{C:flat_nilpotent_kernel}, the
kernel of
\[
 R_k(\mathfrak S_2)\longrightarrow
 R_{k\otimes\mathbb Q}(\mathfrak S_2)=\mathbb Z^2
\]
contains a subgroup isomorphic to
$\bigoplus_{r\geq0}\mathbb Z/2$.
In particular the extra degree-two representation group need not even
be finitely generated, although $k$ is noetherian, local and flat over
the mixed-characteristic DVR and every integer-localized $K_0$ is
$\mathbb Z$.
\end{corollary}
\begin{proof}
For each $r\geq0$ put $d=2^r$ and choose an element
$\alpha\in\mathbb F_{2^d}$ of trace one to $\mathbb F_2$.
Such an element exists: the trace polynomial
$x+x^2+\cdots+x^{2^{d-1}}$ is a nonzero polynomial of degree less
than $2^d$, and its values lie in $\mathbb F_2$.
The element $\alpha$ has degree exactly $d$ over $\mathbb F_2$.
Indeed if it belonged to a proper subfield $\mathbb F_{2^e}$, the
trace would be $d/e$ times its subfield trace, hence zero since $d/e$
is even. Let $q_r(u)$ be its monic irreducible polynomial.
The equation $z^2+z=\alpha$ has no solution in $\mathbb F_{2^d}$:
the map $z\mapsto z^2+z$ has kernel $\mathbb F_2$, and its image is
exactly the trace-zero subspace by dimension and the trace identity.
It follows that
\[
 f_r(t)=q_r(t^2+t)
\]
is irreducible of degree $2d$ over $\mathbb F_2$. More explicitly a
root $\beta$ generates $\alpha=\beta^2+\beta$ and a quadratic
extension of $\mathbb F_{2^d}$, so its degree is $2d$.
The polynomials $f_r$ are pairwise distinct and satisfy
$f_r(t+1)=f_r(t)$.

Lift the entries of a companion matrix of $q_r$ to $0,1$ in the
coefficient DVR, obtaining $B_r$. Its determinant is odd, since
$q_r(0)=1$, so $B_r$ is invertible. Define the free involution
\[
 T_r=\begin{pmatrix}I_d&\epsilon B_r\\
                    \epsilon I_d&-I_d\end{pmatrix},\qquad
 \eta_r=[(k^{2d},T_r)]-d[\mathbf1]-d[\mathrm{sgn}].
                                                               \tag{C.1.8}
\]
Its square is one. Under $\epsilon\mapsto2$ its Bockstein homology
is the whole $\mathbb F_2^{2d}$, with operator
\[
 N_r=\begin{pmatrix}0&\overline B_r\\I_d&I_d\end{pmatrix},
 \qquad \det(tI-N_r)=q_r(t^2+t)=f_r(t).
\]
Consequently $\chi_{f_s}(\eta_r)$ is one for $s=r$ and zero
otherwise. Each class vanishes generically by the two rational
idempotents. The invertible matrix
$J_r=\left(\begin{smallmatrix}0&B_r\\-I_d&0\end{smallmatrix}\right)$
intertwines $T_r$ and $-T_r$, so the evaluation sequence (C.1.7)
gives $2\eta_r=0$. The displayed family of parity homomorphisms
then proves that the subgroup they generate is the asserted direct
sum, without requiring a classification of all involutions.
\end{proof}

The finite certificate \texttt{code/C\_finish\_bockstein.py} constructs
the Bockstein operators from actual kernels and images over
$\mathbb F_2$. It exhausts all $260$ matrices of ranks one and two
over $\mathbb Z/4$, finding $30$ involutions, and all $1808$ block
extension candidates with endpoint ranks $(1,1),(1,2),(2,1)$.
For the $652$ admissible extensions it verifies cyclic homology
exactness, the twisted connecting identity and parity additivity.
It also verifies the rank-four evaluation sequence and the integral
dual-number identities, and the first three independent family members
of dimensions $2,4,8$, whose parity matrix is the identity; output is in
\texttt{logs/C\_finish\_bockstein.json}.

\begin{lemma}[Proved: the dual-number Bockstein detects every nonzero primary factor]
\label{C:epsilon_Bockstein}
Let $A$ be a discrete valuation ring of characteristic zero whose
maximal ideal is $(2)$, put $F=A/2A$, and let
$k=A[\epsilon]/\epsilon^2$. For a finite free $k$-module $M$ with an
involution, write $L=M/\epsilon M$ and
\[
 V(M)=\widehat H^0(C_2,L),\qquad
 W(M)=\widehat H^1(C_2,L).
\]
These are finite-dimensional $F$-spaces. The connecting maps of
$0\to L\xrightarrow{\epsilon}M\to L\to0$, regarded as a sequence of
$A[C_2]$-modules, give natural maps
\[
 b^0:V(M)\longrightarrow W(M),\qquad
 b^1:W(M)\longrightarrow V(M).
\]
Let $u=b^1b^0$ on $V(M)$ and $u=b^0b^1$ on $W(M)$.
For every monic irreducible $q\in F[u]$ other than $u$, the function
\[
 \chi_q^\epsilon(M)
   =\operatorname{length}_{F[u]}V(M)_q\pmod2               \tag{C.1.9}
\]
is additive on admissible short exact sequences and hence defines
$R_k(C_2)\to\mathbb Z/2$.
\end{lemma}
\begin{proof}
Use the two-periodic Tate complex whose alternating differentials
are $T-1$ and $T+1$. Its cohomology is annihilated by $2$, since their
sum is $2$, and is finite over $A$. Thus it is finite-dimensional
over $F$. The sequence defining $b^i$ is exact because $M$ is free
over $k$, and its terms are finite free over $A$.

For clarity, the connecting maps do not necessarily satisfy
$b^1b^0=0$. Choose an $A$-linear section $L\to M$, and write
\[
 T=T_0+\epsilon U,\qquad T_0^2=1,\quad T_0U+UT_0=0.
\]
If $x$ is a Tate cycle in $L$, its lift has differential
$\epsilon Ux$; thus $b^i[x]=[Ux]$. This also follows directly
from the connecting-map definition and makes its independence of
the chosen section explicit.

An admissible sequence $0\to M_A\to M_E\to M_B\to0$ remains
exact after reduction modulo $\epsilon$, because it splits as a
sequence of $k$-modules. Its Tate sequence is the exact six-cycle
\[
 V_A\longrightarrow V_E\longrightarrow V_B
 \xrightarrow{\partial^0}W_A\longrightarrow W_E
 \longrightarrow W_B\xrightarrow{\partial^1}V_A.       \tag{C.1.10}
\]
The Bocksteins commute with the first two maps and anticommute
with the connecting maps. Here is a direct check of the latter
assertion, including its hypotheses. Choose an underlying $k$-linear
splitting and write the two-periodic differential of the middle
complex in block form
\[
 d_E^i=\begin{pmatrix}d_A^i&X\\0&d_B^i\end{pmatrix},
 \quad d_A^i=d_{A,0}^i+\epsilon U_A,\quad
 d_B^i=d_{B,0}^i+\epsilon U_B,\quad X=X_0+\epsilon X_1.
\]
The coefficient of $\epsilon$ in the off-diagonal block of
$d_E^{i+1}d_E^i=0$ is
\[
 U_AX_0+d_{A,0}^{i+1}X_1+X_1d_{B,0}^i+X_0U_B=0.
\]
On a $d_{B,0}^i$-cycle the middle two terms give a boundary.
Since $\partial^i[x]=[X_0x]$, this proves
$b_A^{i+1}\partial^i=-\partial^{i+1}b_B^i$.
In particular $u$ commutes with every map in (C.1.10).

Take the $q$-primary parts. Primary decomposition is exact for
finite-dimensional $F[u]$-modules, so the resulting six-cycle is
exact. Since $q\ne u$, the operator $u$ is invertible on these
parts, and both $b^0$ and $b^1$ are isomorphisms between the
$V$- and $W$-parts. Using $(b_A^0)^{-1}$ to identify $W_{A,q}$
with $V_{A,q}$ turns the first half of (C.1.10) into an exact
three-cycle
\[
 V_{A,q}\longrightarrow V_{E,q}\longrightarrow V_{B,q}
 \xrightarrow{(b_A^0)^{-1}\partial^0}V_{A,q}.
\]
Exactness at its last vertex follows from the commutation of $b^0$
with $V_A\to V_E$ and $W_A\to W_E$. Every map here is
$F[u]$-linear. If its three image lengths are $r,s,t$, the space
lengths are $t+r,r+s,s+t$, respectively. Their alternating sum
is even. This proves (C.1.9).
\end{proof}

\begin{theorem}[Proved: the entire degree-two group over flat dual numbers]
\label{C:dualnumber_C2_ring}
Keep the assumptions of Lemma~\ref{C:epsilon_Bockstein}; in
particular $A=\mathbb Z_{(2)}$ and $A=\mathbb Z_2$ are allowed.
Let $\mathcal I_F$ be the set of monic irreducible polynomials in
$F[u]$ other than $u$. For $q\in\mathcal I_F$ of degree $d$, lift
the entries of its companion matrix to an invertible matrix
$B_q\in M_d(A)$ and set
\[
 T_q=\begin{pmatrix}I_d&\epsilon B_q\\
                     \epsilon I_d&-I_d\end{pmatrix},
 \quad M_q=(k^{2d},T_q),\qquad
 \eta_q=[M_q]-d(1+s),\quad s=[\mathrm{sgn}].             \tag{C.1.11}
\]
Then, as an abelian group,
\[
 R_k(C_2)=\mathbb Z\{1,s\}\ \oplus
             \bigoplus_{q\in\mathcal I_F}(\mathbb Z/2)\eta_q.
                                                               \tag{C.1.12}
\]
For the internal tensor product on this representation group,
the complete multiplication rules are
\[
 s^2=1,\qquad s\eta_q=\eta_q,\qquad
 \eta_q\eta_{q'}=0\quad(q,q'\in\mathcal I_F).           \tag{C.1.13}
\]
The second summand is precisely the kernel of generic scalar
extension to $R_{k[1/2]}(C_2)=\mathbb Z\{1,s\}$ and also
the cokernel of the degree-two polynomial induction classes.
Thus this is a computation of the full degree-two obstruction,
not merely an injection into that obstruction. The assertion
concerns this family of base rings; it does not settle C.1 for
every commutative ring.
\end{theorem}
\begin{proof}
\emph{The integral starting point.}
We first give an elementary decomposition of every finite free
$A$-module $L$ with an involution $T_0$. Put
$L_\pm=\ker(T_0\mp1)$. These are finite free over the DVR $A$,
their direct sum has full rank in $L$, and $2L\subset L_++L_-$.
The quotient $Q=L/(L_++L_-)$ is an $F$-space, and the maps
\[
 Q\longrightarrow L_+/2L_+,\ [v]\longmapsto[v+T_0v],
 \qquad
 Q\longrightarrow L_-/2L_-,\ [v]\longmapsto[v-T_0v]
\]
are injective. For example if $v+T_0v=2a$ with $a\in L_+$,
then $v-a\in L_-$; the other assertion is identical.
Choose an $F$-basis $[v_1],\ldots,[v_r]$ of $Q$.
The elements $e_i=v_i+T_0v_i$ extend to an $A$-basis of $L_+$,
and the elements $f_i=v_i-T_0v_i$ extend to one of $L_-$,
because their reductions are independent. The relations
$2v_i=e_i+f_i$ show that $L$ has an $A$-basis consisting of
$v_i,T_0v_i$ for $1\le i\le r$, followed by the unused basis
vectors of $L_+$ and $L_-$. On each pair $T_0$ swaps the two
vectors. We have proved
\[
 L\simeq A[C_2]^r\oplus A^a\oplus A_{\mathrm{sgn}}^b.  \tag{C.1.14}
\]

\emph{Lifting and removing the regular summands.}
For an arbitrary object $M$ over $k$, apply (C.1.14) to
$L=M/\epsilon M$. Lift each regular generator $v_i$ to $M$.
The vectors $v_i,Tv_i$, together with lifts of the remaining
basis of $L$, form a $k$-basis of $M$: their change-of-basis
matrix is invertible modulo the nilpotent ideal $(\epsilon)$.
They therefore define a $k[C_2]^r$-submodule with a free
underlying quotient. It splits equivariantly. Indeed, choose
$k$-linear functionals $a_i$ taking values $1,0$ on
$v_i,Tv_i$ and zero on the other basis vectors. The maps
\[
 m\longmapsto a_i(m)+g\,a_i(Tm)\in k[C_2]
\]
are equivariant and jointly retract this submodule. Modulo
$\epsilon$ this is the original regular projection, so the
complement has reduction $A^a\oplus A_{\mathrm{sgn}}^b$.
Regular summands contribute only $1+s$ to $R_k(C_2)$.

\emph{A finite-length matrix problem gives all generators.}
On the complement, lift a basis of its plus and minus reduction.
Its involution has the form
\[
 T=\begin{pmatrix}I_a&0\\0&-I_b\end{pmatrix}
    +\epsilon\begin{pmatrix}0&B\\ C&0\end{pmatrix}.    \tag{C.1.15}
\]
Indeed $T^2=1$ forces the two diagonal blocks of the first-order
term to be zero, since $A$ has no $2$-torsion.
Conjugation by $I+\epsilon X$ changes $B$ and $C$ independently
by arbitrary matrices divisible by $2$. Consequently only
their reductions
\[
 \overline C:F^a\longrightarrow F^b,\qquad
 \overline B:F^b\longrightarrow F^a
\]
matter for the following filtration argument. An invertible
change of basis at the two $F$-spaces lifts to one over $A$.
For a subrepresentation of this two-arrow cycle choose bases
adapted to its two subspaces. The matrix blocks which must vanish
for invariance are divisible by $2$; an $I+\epsilon X$
conjugation makes them exactly zero. The lifted coordinate
subspaces therefore give an invariant free $k$-submodule with
free quotient. Repeating along a composition series expresses
$[M]$ as a sum of the lifts of the simple representations of
this finite-dimensional quiver representation.

Its simple representations can be classified directly. There
are the two one-dimensional vertex representations. If both
vertex spaces of a simple representation are nonzero, the
kernel of either arrow would give a subrepresentation supported
at one vertex; hence both arrows are injective and thus are
isomorphisms. Identify one arrow with the identity. Simplicity
then says that the other is a simple invertible $F[u]$-module,
so it is the companion matrix of one uniquely determined
$q\in\mathcal I_F$. Lifting its basis change and adjusting even
entries in (C.1.15) gives precisely $M_q$. We have proved that
$1,s$ and the $[M_q]$, hence $1,s,\eta_q$, generate the group.
This uses finite-length filtrations to compute the group; it
does not claim that all indecomposable dual-number lattices
have been classified.

\emph{All relations and the generic kernel.}
The determinant of $B_q$ is a unit since $q(0)\ne0$.
The invertible matrix
\[
 J_q=\begin{pmatrix}0&B_q\\-I_d&0\end{pmatrix}
\]
satisfies $J_qT_q=(-T_q)J_q$. Thus $s[M_q]=[M_q]$.
The evaluation sequence (C.1.7), of rank $2d$ here, gives
$2[M_q]=2d(1+s)$, and therefore $2\eta_q=0$ and
$s\eta_q=\eta_q$.

The reduction of $M_q$ modulo $\epsilon$ is the direct sum of
$d$ plus and $d$ minus lines. Its Tate groups are respectively
$F^d$ and $F^d$, and its Bocksteins are
\[
 b^0=I_d,\qquad b^1=\overline B_q.
\]
It follows that
$\chi_{q'}^\epsilon(\eta_q)=1$ if $q'=q$ and is zero otherwise.
The two constant lines have zero Bockstein, so every such
invariant vanishes on them. Generically $(1\pm T_q)/2$ split
$M_q$ into free plus and minus modules of rank $d$: the ring
$k[1/2]=\operatorname{Frac}(A)[\epsilon]/\epsilon^2$ is local.
Thus the generic comparison records the two independent integer
coefficients of $1,s$, and the homomorphisms
$\chi_q^\epsilon$ record independently every proposed
$\mathbb Z/2$ coefficient. These detectors, together with the
generation already proved, give exactly (C.1.12) and its asserted
generic kernel. The degree-two induction image is the span of
$1$ and the regular class $1+s$, which is the same integer
summand.

\emph{Multiplication.}
Let $B=B_q$ and $C=B_{q'}$ have sizes $d,e$. Reduce
$M_q\otimes_k M_{q'}$ modulo $\epsilon$. Its plus Tate space
is the sum of the plus--plus and minus--minus tensor spaces,
each of dimension $de$. The first-order involution is
$U_B\otimes T_{C,0}+T_{B,0}\otimes U_C$. Modulo $2$ its square
on the plus Tate space is the direct sum of two copies of
\[
 \overline B\otimes I_e+I_d\otimes\overline C.          \tag{C.1.16}
\]
This follows by squaring: the two cross terms cancel in
characteristic two and $U_B^2$ acts by $B$ on each sign space,
and similarly for $C$. Hence every primary length is even and
every $\chi_{q''}^\epsilon$ of this tensor product is zero.
Its generic class is $de(1+s)^2=2de(1+s)$; by (C.1.12),
\[
 [M_q][M_{q'}]=2de(1+s).
\]
Since $(1+s)\eta_q=2\eta_q=0$, substituting
$[M_q]=d(1+s)+\eta_q$ gives $\eta_q\eta_{q'}=0$.
This proves all of (C.1.13).
\end{proof}

For $A=\mathbb Z_{(2)}$ or $\mathbb Z_2$, the index set in
(C.1.12) is countably infinite. The earlier map to $\mathbb Z/4$
detects those $q$ for which $q(t^2+t)$ is irreducible; the
$\epsilon$-Bockstein also detects the other irreducibles.
For example $q=u^3+u+1$ has trace-zero roots and
$q(t^2+t)$ splits into two cubics exchanged by $t\mapsto t+1$,
whereas $\chi_q^\epsilon(\eta_q)=1$. This explains why the
first Bockstein calculation alone could only prove a subgroup
injection and why the full calculation requires another invariant.

The finite certificate \texttt{code/C\_finish\_dualnumber\_ring.py}
constructs the integral Tate groups from actual kernels, images and
the modulo-four Bockstein, then computes the $\epsilon$-Bockstein
from the first-order matrix. Among $6736$ block extension candidates
on seven endpoint objects it checks all $197$ admissible extensions
for cyclic exactness, Bockstein compatibility and primary parity
additivity. It checks all $292$ two-cycle matrix pairs over
$\mathbb F_2$ in dimensions $(1,1),(1,2),(2,1),(2,2)$, lifting
all $2076$ invariant subspace pairs by integral even-entry gauges.
Finally it checks the seven nonzero irreducibles of degrees at most
four and all $49$ corresponding tensor identities (C.1.16).
The finite scope, matrices and runtime are recorded in
\texttt{logs/C\_finish\_dualnumber\_ring.json}; the general
classification rests on the proof above.

\begin{proposition}[Proved: nil-invariance in uniquely divisible directions]
\label{C:divisible_nil_invariance}
Let $k$ be any commutative ring, $G$ a finite group of order $N$,
and $J\subset k$ a nil ideal such that multiplication by $N$ is a
bijection on the additive group of $J$. Then reduction induces
\[
 R_k(G)\xrightarrow{\ \sim\ }R_{k/J}(G).                \tag{C.1.17}
\]
It respects internal tensor products and exterior operations.
In particular, if the additive group of $J$ is a $\mathbb Q$-vector
space, reduction gives an isomorphism of the entire induction rings
$R_k(\mathfrak S_*)\simeq R_{k/J}(\mathfrak S_*)$, and
$\vartheta_k$ is an isomorphism if and only if $\vartheta_{k/J}$ is.
There is no noetherian or finite-generation hypothesis on $J$.
\end{proposition}
\begin{proof}
First suppose $J^2=0$. Finite projectives lift across $J$: represent
one by a finite idempotent matrix, lift the matrix, and replace an
approximate idempotent $e$ by $3e^2-2e^3$; its error is a multiple
of $(e^2-e)^2$ and hence zero. This is the finite case of
\cite[Tags 07LV and 0AMF, Lemma 10.32.7]{C-StacksNil}.
Maps between such projectives lift by projectivity, and any lift
of an isomorphism is invertible, by a finite geometric series in
the square-zero error.

Fix a projective lift $P$ of a representation $\overline P$ over
$k/J$. Lift each of the finitely many action matrices to an
automorphism $r_g$ of $P$, taking $r_1=1$. Let
$D=\operatorname{Hom}_k(P,JP)$; it is square-zero as an ideal in
$\operatorname{End}_k(P)$, and $N$ acts invertibly on $D$ because
$P$ is a direct summand of a finite free module. Conjugation by
the $r_g$ defines an actual $G$-action on $D$: all group-law
errors lie in $D$, whose conjugation action on itself is trivial.
Write
\[
 r_gr_hr_{gh}^{-1}=1+c(g,h),\qquad c(g,h)\in D.
\]
Associativity gives
\[
 c(g,h)+c(gh,l)=g\,c(h,l)+c(g,hl).
\]
Define $a(g)=N^{-1}\sum_{l\in G}c(g,l)$, where division takes
place in $D$; it does not require $N$ to be a unit in $k$.
Summing the preceding identity over $l$ gives
\[
 c(g,h)=a(g)+g\,a(h)-a(gh).
\]
Consequently $r'_g=(1-a(g))r_g$ obeys the group law. Every
representation over $k/J$ has therefore lifted to one over $k$.

Next let $P,Q$ already carry lifted actions, and let
$\bar f:\overline P\to\overline Q$ be equivariant. Choose an
underlying $k$-linear lift $f$. The discrepancies
\[
 d(g)=g\,f-f\in\operatorname{Hom}_k(P,JQ)
\]
obey $d(gh)=d(g)+g\,d(h)$. For
$b=N^{-1}\sum_h d(h)$, summing this identity gives
$d(g)=b-g\,b$. Hence $f+b$ is an equivariant lift. In particular
an equivariant isomorphism lifts to one. This proves uniqueness
of lifts up to isomorphism as well as surjectivity on all Hom sets;
it does not assert that reduction is faithful on Hom sets.

Admissible relations lift too. Given
$0\to\overline A\to\overline E\to\overline B\to0$, lift its
three objects and lift the equivariant surjection to $p:E\to B$.
An underlying $k$-linear lift $s:B\to E$ of its chosen splitting
has $ps$ invertible, so replacing $s$ by $s(ps)^{-1}$ makes
$p$ a split surjection. Its kernel is projective over $k$ and
reduces to $\overline A$; the uniqueness just proved identifies
it equivariantly with the chosen lift $A$. Thus both the objects
and all defining relations of the two Grothendieck groups
correspond, proving (C.1.17) when $J^2=0$.

If $J$ is nilpotent, apply this argument successively to
$J^r/J^{r+1}$. Multiplication by $N$ is bijective on every power
$J^r$: surjectivity follows by dividing one factor of each product,
and injectivity follows from $J^r\subset J$. It is consequently
bijective on each quotient layer. The square-zero argument applies
at every stage.

Finally let $J$ merely be nil. The inverse of multiplication by
$N$ is $k$-linear on $J$. For a finite set $j_1,\ldots,j_t\in J$,
let $L$ be the ideal it generates, and put
\[
 L^{[N]}=\sum_{a\geq0}N^{-a}L\ \subset J.
\]
This is an ideal on which $N$ acts bijectively. It is nilpotent:
the finitely generated nil ideal $L$ has some $L^r=0$, and a
product of $r$ elements of $L^{[N]}$, after multiplication by a
suitable power of $N$, is a sum of products from $L^r$.
Injectivity of that power on $J$ forces the product to be zero.
The ideals $L^{[N]}$ form a filtered family with union $J$.
Thus $k/J=\operatorname*{colim}_L k/L^{[N]}$.
Proposition~\ref{C:filtered_colimits}, whose finite-matrix descent
also includes equivariant maps and admissible sequences, now
extends (C.1.17) from every $L^{[N]}$ to $J$.

Reduction commutes with tensor products, exterior powers and
induction. If $J$ is a $\mathbb Q$-vector space the hypothesis
holds for every $N$, so these degreewise isomorphisms give the
last assertion and carry the polynomial generators to themselves.
\end{proof}

\begin{corollary}[Proved: nonreduced mixed-characteristic positive examples]
\label{C:divisible_nil_positive}
Let $A$ be any ring satisfying C.1 and let $J$ be a nil ideal in
a commutative extension $k$ with $k/J\simeq A$. If $J$ is uniquely
divisible as an additive group, then $k$ satisfies C.1.
For example
\[
 k=\mathbb Z_{(2)}\oplus\mathbb Q\epsilon,
 \qquad\epsilon^2=0,\qquad
 (a,x)(b,y)=(ab,ay+bx)                                \tag{C.1.18}
\]
is local, nonreduced and nonnoetherian, and
$\vartheta_k$ is an isomorphism in every degree. Its closed
fibre is $\mathbb F_2$, its generic fibre is
$\mathbb Q[\epsilon]/\epsilon^2$, and $K_0(k)=\mathbb Z$.
\end{corollary}
\begin{proof}
Apply Proposition~\ref{C:divisible_nil_invariance}.
In (C.1.18) the maximal ideal is
$2\mathbb Z_{(2)}\oplus\mathbb Q\epsilon$; reduction by the
square-zero rational ideal is $\mathbb Z_{(2)}$, covered by
Proposition~\ref{C:DVR}. The ideal $\mathbb Q\epsilon$ is not
finitely generated over $k$: a finite collection has bounded
$2$-denominators, and multiplying by $k$ cannot enlarge them.
This proves nonnoetherianity. The fibre statements follow by
tensoring with $\mathbb F_2$ and $\mathbb Q$, respectively, and
finite projectives over the local ring are free.
\end{proof}

\begin{corollary}[Proved: ordinary fibre and projective-class data are insufficient]
\label{C:fibre_invisible_kernel}
For $A=\mathbb Z_{(2)}$ the maps
$A\to A[\epsilon]/\epsilon^2\to A$ induce the same underlying
spectrum and residue fields. On every integer localization, and
on every base change to a field, their ordinary $K_0$ groups are $\mathbb Z$
and these maps induce the identity. Nevertheless the first ring
satisfies C.1 and the second has the infinite kernel (C.1.12).
Every element of that kernel is zero on restriction to each
residue-field point of $\operatorname{Spec}k$. Thus ordinary
$K_0$ data and these pointwise representation classes alone cannot
give a necessary and sufficient test for C.1. This does not rule
out a criterion that also uses further structure of the ring $k$.
\end{corollary}
\begin{proof}
The underlying spectrum is unchanged by a nilpotent ideal.
The local rings in question, their integer localizations, and
their base changes to any residue field are local; hence their
$K_0$ groups are $\mathbb Z$, with the unit classes matched.
The positive assertion for $A$ is Proposition~\ref{C:DVR}, while
Theorem~\ref{C:dualnumber_C2_ring} gives the full kernel for the
dual-number ring. At its characteristic-zero point each $\eta_q$
vanishes generically. At its characteristic-two point
$\epsilon=0$ and both signs become trivial, so the two terms
of (C.1.11) are the same $2\deg(q)$ trivial representation.
This proves vanishing at every field point without discarding
the nonzero integral classes.
\end{proof}

The finite certificate \texttt{code/C\_finish\_nilfibres.py}
checks the two averaging homotopy identities before division by
the group order on every basis cochain for $C_2,C_3,\mathfrak S_3$
acting on the endomorphisms of their indicated permutation modules:
$89$ basis vectors in degree one and $421$ in degree two. These
integer identities verify the signs used to lift actions and maps;
the arbitrary-ring extension is the proof above.


\begin{remark}[Open: the unsolved all-rings classification]
\label{C:allrings_open}\label{C:allrings_gate}
Equation (C.1.1) is an exact obstruction reduction, not a claim to have
classified every commutative ring. The results above classify all
$\mathbb Q$-algebras and all regular noetherian rings, exclude every
positive-characteristic ring, and establish stability under filtered
colimits. Theorem~\ref{C:regular_modular_locus} further allows every
characteristic-zero singularity and requires regularity only on the
positive-characteristic locus. They also disprove the proposed criteria of locality,
$K_0(k)=\mathbb Z$, connected regular rational fibre, and even
$K_0(k[1/N])=\mathbb Z$ for every integer $N$, in the stated examples.
What remains is to calculate both groups in (C.1.1) for arbitrary
mixed-characteristic rings with singularities on that positive-characteristic
locus and nonnoetherian bases beyond the filtered-colimit cases.
Proposition~\ref{C:flat_nilpotent_kernel} now settles the previously
undecided flat dual-number ring over $\mathbb Z_{(2)}$ negatively,
and proves that even flatness does not remove the generic kernel.
Theorem~\ref{C:dualnumber_C2_ring} calculates that entire degree-two
kernel over any characteristic-zero DVR with maximal ideal $(2)$:
there is one independent order-two class for each nonzero irreducible
polynomial over the residue field, and their internal tensor products
are zero. This resolves this family completely in degree two but
does not supply the missing criterion for arbitrary singular bases.
Proposition~\ref{C:divisible_nil_invariance} removes nil ideals
that are uniquely divisible by all integers, without any noetherian
hypothesis, and gives nonreduced mixed-characteristic positive examples.
The remaining gap concerns integral nil directions as well as reduced
singular bases; unrestricted nilpotent invariance and generic-fibre
injectivity are both false.
\end{remark}

\subsection{C.2: the ring actually defined by the finite-set pages}

Let $\mathrm{FI}$ be the category of finite sets and injections. A finite
species is a functor from finite sets and bijections to finite sets, with
$X[I]=\varnothing$ for $|I|$ sufficiently large. Write
\[
 \mathcal F_X(E)=\coprod_{A\subseteq E}X[A]
 \simeq\coprod_{n\geq0}\operatorname{Inj}([n],E)
                 \mathbin{\times}_{\mathfrak S_n}X[n].
\]
The second expression uses the usual balanced quotient, with the action
conventions chosen accordingly. This is exactly the finite-support reading
of the blue marginal formula on cote 133, scan 26.

\begin{theorem}[Proved: distinguished functors and eventual equivalence]
\label{C:distinguished_ring}
The following give the complete additive and multiplicative description of
the Grothendieck ring $\mathcal D$ of distinguished functors modulo the
manuscript's eventual equivalences.
\begin{enumerate}
\item Each distinguished functor is uniquely of the form $\mathcal F_X$ up
to isomorphism of finite species. An eventually bijective natural
transformation between distinguished functors is bijective at every finite
set. Even if zigzags through arbitrary functors are allowed, eventual
equivalence identifies no two nonisomorphic distinguished functors.
\item As an abelian group,
\[
 \mathcal D=\bigoplus_{n\geq0}B(\mathfrak S_n)
 =\bigoplus_{n\geq0}\ 
   \bigoplus_{[H]\subseteq\mathfrak S_n}\mathbb Z\,[\operatorname{Inj}([n],-)/H],
\]
where $[H]$ runs through conjugacy classes of subgroups.
\item Pointwise product is the union-of-supports product
\[
 (X\star Y)[I]=\coprod_{A,B\subseteq I,\ A\cup B=I}X[A]\times Y[B].
\]
In particular this is not the usual species product that requires
$A\cap B=\varnothing$. The unit is the constant singleton functor.
\item Composition on effective functors is given explicitly by taking,
among the elements of $\mathcal F_X(\mathcal F_Y(I))$, those for which the
union of the supports in $I$ of all selected $Y$-structures is $I$.
This gives a finite species of degree at most $\deg(X)\deg(Y)$.
Composition extends uniquely to virtual classes, additively in its outer
argument and polynomially in its inner argument, and remains associative.
Its unit is $E\mapsto E$, not the constant singleton.
\end{enumerate}
\end{theorem}
\begin{proof}
An element $(A,x)$ of $\mathcal F_X(E)$ has the unique least support $A$:
it lies in the image from $\mathcal F_X(B)$ exactly when $A\subseteq B$.
The elements of full support in $\mathcal F_X(I)$ therefore recover $X[I]$,
functorially for bijections. Each $X[n]$ decomposes uniquely into transitive
$\mathfrak S_n$-sets. These observations prove both uniqueness and the
displayed additive basis.

For an eventual isomorphism on the full tail of $\mathrm{FI}$, choose a set
$E$ larger than the tail bound and with $|E\setminus A|>\deg(Y)$ for a
prescribed support $A$. The stabilizer of an element supported in $A$
contains every permutation of $E\setminus A$. Equivariance forces the
support of its image to be invariant under these permutations. If that
support met $E\setminus A$, it would contain the entire complement,
contradicting its degree bound. Thus its support is contained in $A$;
this argument applies even when the natural isomorphism is defined only
on the tail and $A$ itself is not in its domain. If a bijection on the
tail decreased the support of an element from $A$ to $B\subsetneq A$, a
transposition fixing $B$ and exchanging an element of $A\setminus B$ with
an element outside $A$ would fix its image but not the original element.
This contradicts injectivity. Hence tail isomorphisms preserve supports.
Their bijections on the elements supported in $A$ recover isomorphisms
$\mathcal F_X(A)\simeq\mathcal F_Y(A)$ for every $A$. They are independent
of the chosen large ambient set, by embedding any two choices in a larger
one and using naturality there. A zigzag of eventual
isomorphisms restricts to an isomorphism on a sufficiently large common
tail, so the same argument applies.

A pair of elements with supports $A$ and $B$ has support $A\cup B$.
This proves the product formula, with its actual $\mathfrak S_I$-action,
and proves that it lists every product summand. An element of a composite
selects at most $\deg(X)$ distinct $Y$-structures. The union of their
supports has at most $\deg(X)\deg(Y)$ elements. This proves the stated
composition formula and finiteness, including constant functors.

For completeness, fix an outer orbit functor
$F=\operatorname{Inj}([n],-)/H$. Its $(n+1)$st cross-effect for disjoint
union vanishes: a selection of $n$ structures cannot meet $n+1$ prescribed
nonempty colour classes. Thus, on the free commutative monoid of
distinguished isomorphism classes, $G\mapsto[F\circ G]$ is polynomial of
degree at most $n$. On every finite collection of monoid generators,
successive finite differences express it uniquely as an integral linear
combination of products of binomial polynomials in their multiplicities.
These polynomials have integer values on all integers, giving a unique
extension to the group completion. The constructions agree upon adding
unused generators, so define the extension on $\mathcal D$.
Associativity and the identities involving outer addition and product hold
on nonnegative multiplicities. Both sides are polynomial on every finite
collection of generators, hence hold on all integer multiplicities by
the same finite-difference uniqueness. This proves the last assertion.
\end{proof}

\begin{proposition}[Refuted: L7 as a special-$\lambda$ claim]
\label{C:lambda_failure}
The finite-set constructions produce two different natural pre-$\lambda$
structures on $\mathcal D$, neither of which is special:
\[
\begin{aligned}
 a^n(F)(E)&=\{n\text{-element subsets of }F(E)\},\\
 \lambda_t(F)&=\sigma_{-t}(F)^{-1},\qquad
 \sigma^n(F)(E)=\operatorname{Sym}^n(F(E)).
\end{aligned}
\]
In particular the subset powers in scan 19 cannot simultaneously satisfy
the inverse-symmetric-power identity written there. The corrected structure
retains all the subgroup-indexed operations and composition in
Theorem~\ref{C:distinguished_ring}; discarding them in favour of a special
$\lambda$-ring loses information.
\end{proposition}
\begin{proof}
Decomposing a subset or multiset according to two disjoint summands proves
the pre-$\lambda$ addition identity. Both constructions preserve
distinguished functors by the support description, so they extend to
$\mathcal D$ by group completion.

Let $C_2$ act freely on a two-point set $X$ and write $x=[X]$ in its
Burnside ring. Then $x^2=2x$ and $a^2(x)=1$. The two-subsets of $X^2$
have two fixed points and two free orbits, so
\[
 a^2(x^2)=2+2x\ne4x-2
 =x^2a^2(x)+x^2a^2(x)-2a^2(x)^2.
\]
The full-group fixed-point marks of the two sides are $2$ and $-2$.
Also $\sigma^2(x)=x+1$, whereas the inverse-symmetric rule gives
$\lambda^2(x)=x^2-\sigma^2(x)=x-1\ne1=a^2(x)$.

For the inverse-symmetric structure let $V_4$ act regularly on a four-point
set $E$ and put $X=\mathcal P_2(E)$. Its six points are the disjoint union
of the three transitive two-point $V_4$-sets. Thus $X$ has no fixed points,
$\operatorname{Sym}^2(X)$ has three, and
$\operatorname{Sym}^2(X\times X)$ has six. If $m$ denotes the full-group
fixed-point mark, we get
\[
 m(\lambda^2[X])=-3,\qquad
 m(\lambda^2([X]^2))=-6,
\]
while the mark of the special-$\lambda$ product polynomial is
$-2(-3)^2=-18$. The defect is $12$.
These are evaluations of distinguished functors: use $F(E)=E$ in the
first calculation and $F(E)=\mathcal P_2(E)$ in the second. Evaluation at
a finite set with a group action is compatible with sum, product, and all
the operations used. Theorem~\ref{C:distinguished_ring} shows that eventual
equivalence cannot remove the defect. Direct enumeration of all the sets
and marks is included in \texttt{code/C\_finite\_certificates.py}.
The inverse-symmetric convention agrees with Bouc--R\"okaeus,
Introduction and Theorem 1.1 \cite{C-BoucRokaeus2009}; the subset convention
does not.
\end{proof}

\begin{proposition}[Proved: the corrected partition-type formula]
\label{C:partition_formula}
For any group $G$ and finite $G$-set $E$, let
$\tau^m(E)=[\operatorname{Inj}([m],E)]\in B(G\times\mathfrak S_m)$ and
$T^n(E)=[E^{[n]}]\in B(G\times\mathfrak S_n)$. For a partition type
$\alpha=(\alpha_1,\ldots,\alpha_n)$ with $\sum i\alpha_i=n$, put
$m=\sum\alpha_i$, choose a partition of $[n]$ with this type, and let
\[
 W_\alpha=\prod_i(\mathfrak S_i\wr\mathfrak S_{\alpha_i}),\qquad
 q_\alpha:W_\alpha\longrightarrow\prod_i\mathfrak S_{\alpha_i}
                               \hookrightarrow\mathfrak S_m
\]
be its action on the set of blocks. Define the additive homomorphism
\[
 \phi_\alpha=\operatorname{Ind}_{G\times W_\alpha}^{G\times\mathfrak S_n}
       \circ(\operatorname{id}_G\times q_\alpha)^*.
\]
Then, including $n=0$ with the empty partition convention,
\[
 T^n(x)=\sum_{\alpha:\,\sum i\alpha_i=n}\phi_\alpha(\tau^{\sum\alpha_i}(x))
 \quad\text{for every }x\in B(G).
\]
The pullback in $\phi_\alpha$ is restriction to
$\prod_i\mathfrak S_{\alpha_i}$ followed by inflation; inflation directly
from all of $\mathfrak S_m$ would be an inaccurate description unless this
restriction is understood.
\end{proposition}
\begin{proof}
The kernel partition of a function $[n]\to E$ is preserved by $G$ and
permuted by $\mathfrak S_n$. The partitions of type $\alpha$ form a
transitive $\mathfrak S_n$-set with stabilizer $W_\alpha$. Over the chosen
partition, a function is exactly an injection from its $m$ blocks to $E$,
and the stabilizer acts through $q_\alpha$. This gives the displayed
isomorphism of actual $G\times\mathfrak S_n$-sets, after taking the disjoint
union over all types. For virtual $G$-sets, $T^n$ and each $\tau^m$ have
polynomial degree at most $n$, by the same colour-class cross-effect
argument as above. Equality on actual finite $G$-sets therefore extends
uniquely to the group completion. Finally,
$q_H^n(x)$ in scans 21--23 is the additive orbit map
$B(G\times\mathfrak S_n)\to B(G)$ applied to $\tau^n(x)$;
it is a polynomial operation of degree at most $n$, generally not
homogeneous under multiplication of the argument by an integer.
\end{proof}

\subsection{C.3: the split quadratic-form ring}

Throughout this subsection $2$ is invertible. We write $WG_{\mathrm{sp}}(S)$
for the group completion of isometry classes of nondegenerate symmetric
bundles under orthogonal sum. This is scan 6's $WG(S)$, using the equivalence
between quadratic and symmetric bilinear forms when $2$ is invertible.
Exterior powers define a pre-$\lambda$-structure: the decomposition of an
exterior power of an orthogonal sum is an actual orthogonal decomposition.
The word ``special'' below means the additional universal identities for
products and iterated exterior powers. Modern exact Grothendieck--Witt groups
impose, in addition, the relation that a metabolic object with Lagrangian $L$
equals $H(L)$. That relation is not present in the manuscript's definition.

\begin{proposition}[Proved: the characteristic-two boundary of the manuscript's definition]
\label{C:char2_boundary}
For strictly nondegenerate quadratic forms, exterior powers cannot define
the indicated pre-$\lambda$ structure when a characteristic-two fibre is
present. Moreover the ordinary rule
$q_{V\otimes W}(v\otimes w)=q_V(v)q_W(w)$ does not define a functorial
quadratic tensor product even over $\mathbb F_2$. Thus the restriction
$2\in\mathcal O_S^\times$ in the ring statements below is essential.
This concerns quadratic forms, not symmetric bilinear forms, which are
different objects in characteristic two.
\end{proposition}
\begin{proof}
Over a field of characteristic two the polar form of a quadratic form
is alternating, so strict nondegeneracy forces even dimension.
The hyperbolic quadratic plane is strictly nondegenerate, but its
second exterior power has underlying dimension one. Even a virtual
class of strictly nondegenerate quadratic forms has even rank, so it
cannot supply the required exterior-square class of rank one.
Restriction to a characteristic-two fibre gives the same obstruction
on a base where two is not a unit.

For the tensor assertion take $H=\mathbb F_2^2$ with $q(x,y)=xy$ and
use coordinates $(a,b,c,d)$ on $H\otimes H$ arranged in a two-by-two
matrix. The required values on pure tensors force the square terms
and the coefficients of $ab,ac,bd,cd$ to be zero. Only $ad$ and $bc$
remain, and their coefficients must sum to one by evaluation at the
all-ones pure tensor. Interchanging the two coordinates of the first
factor is an isometry of $H$ and exchanges $ad$ with $bc$. Naturality
would make their coefficients equal, a contradiction.
The attached exact enumeration
\texttt{code/C\_characteristic2\_quadratic.py} checks all 1024 quadratic
forms on the tensor square: precisely two have the required pure-tensor
values, and neither is invariant under the factor isometries.
This verifies the restriction already recorded on cote 133, scan 6.
\end{proof}

\begin{proposition}[Proved: L8 in characteristic zero]
\label{C:charzero_lambda}
For every scheme $S$ over a field of characteristic zero,
$WG_{\mathrm{sp}}(S)$ is a special $\lambda$-ring under exterior powers.
In particular L8 holds, with the split convention required by scan 6.
\end{proposition}
\begin{proof}
Let $k$ be the ground field. Work first on the open and closed strata of
$S$ where the ranks of the bundles involved are fixed. The representation categories of $O_m$ and
$O_m\times O_n$ over $k$ are semisimple: their identity components are
reductive in characteristic zero and their component groups have order
invertible in $k$. In a semisimple category with duality and with $2$
invertible, a metabolic symmetric object is isometric to the hyperbolic
object on its Lagrangian. Indeed split the underlying sequence
$0\to L\to V\to L^*\to0$ equivariantly. In this decomposition the form has
matrix $\left(\begin{smallmatrix}0&1\\1&a\end{smallmatrix}\right)$; replacing
the chosen complement by its graph with component $-a/2$ makes the lower
right block zero. Thus the extra metabolic relations are already isometries,
and split $WG$ equals exact $GW$ for these universal representation categories.

Zibrowius's main theorem \cite{C-Zibrowius2015} makes the latter groups
special $\lambda$-rings; see also his Corollary 3.18 for these particular
groups. For symmetric bundles $V,W$ on $S$, the product of their isometry
torsors gives an exact tensor functor from representations of
$O_m\times O_n$ to bundles on $S$. It preserves forms, actual orthogonal
sums, and exterior powers. Applying it to the universal product identity
proves that identity for $V,W$. The isometry torsor of $V$ similarly transfers
the universal iterated-exterior-power identities from $O_m$. These torsors
exist for the etale topology because every nondegenerate symmetric form is
etale locally isometric to the standard form. The identities for differences
of actual bundles follow from the pre-$\lambda$ addition formula, or,
equivalently, by expanding their universal integral polynomials after adjoining
formal additive inverses. Hence they hold on the group completion.
The identities on the rank strata glue: the bundles, forms, isometries,
and any stabilizing orthogonal summands can all be glued over a disjoint
open cover. This also covers schemes with unbounded locally finite ranks.

This argument uses semisimplicity exactly where it identifies split and exact
$GW$. Zibrowius's Definition 1.2 explicitly includes metabolic relations;
his positive-characteristic theorem therefore cannot by itself answer the
split question below.
\end{proof}

\begin{proposition}[Proved: $WG(B_{O_m})$ for split $O_m$ in characteristic zero]
\label{C:orthogonal_representation_ring}
Let $k$ be a field of characteristic zero and let $O_m$ be the split
orthogonal group. Choose one nondegenerate invariant symmetric form $q_\lambda$
on each simple representation $V_\lambda$. Then
\[
 WG_{\mathrm{sp}}(B_{O_m})
   =\bigoplus_{\lambda\in\operatorname{Irr}(O_m)}GW(k)[V_\lambda,q_\lambda]
\]
as a $GW(k)$-module. This gives a complete ring description by the following
finite multiplicity-form rule: the coefficient of $[V_\nu,q_\nu]$ in a
product is the class of the symmetric multiplicity space
\[
 \operatorname{Hom}_{O_m}(V_\nu,V_\lambda\otimes V_\mu),
\]
whose form is uniquely specified by requiring evaluation to identify the
$\nu$-isotypic summand with $V_\nu$ tensored with that space. Exterior powers
are calculated by the identical rule with $\Lambda^r(V_\lambda\otimes U)$.
Over an algebraically closed $k$, forgetting forms identifies this ring and
its $\lambda$-structure with the ordinary representation ring $R(O_m)$.
\end{proposition}
\begin{proof}
Every simple representation of split $O_m$ has endomorphism ring $k$ and
admits a symmetric form: this is Zibrowius, Proposition 3.17(c). In the
semisimple category, distinct simple isotypic components of a symmetric
object are orthogonal because a pairing between them is an equivariant map
to the dual, and each simple is self-dual. On the $\lambda$-component,
Schur's lemma writes the form uniquely as $q_\lambda\otimes b_\lambda$.
It is nondegenerate exactly when $b_\lambda$ is. This proves the displayed
additive decomposition by actual isometries and proves the product and
exterior-power rules. If $k$ is algebraically closed, symmetric forms are
classified by dimension and $GW(k)=\mathbb Z$, giving the last assertion.

For explicit indexing, when $m=2r+1$ the labels are the integral dominant
weights $\lambda_1\geq\cdots\geq\lambda_r\geq0$, with two extensions for each,
interchanged by determinant. When $m=2r$, a weight with
$\lambda_1\geq\cdots\geq\lambda_r\geq0$ and $\lambda_r=0$ has two extensions;
a weight with $\lambda_r>0$ gives one $O_{2r}$-simple whose restriction to
$SO_{2r}$ is the sum of the weight $\lambda$ and the weight obtained by
negating $\lambda_r$. These are precisely Proposition 3.17(a),(b).
The case $O_0=1$ is $GW(k)$. Over a general field, replacing a multiplicity
form by its dimension would lose information, so the coefficient ring
$GW(k)$ and the stated forms are essential.
\end{proof}

\begin{theorem}[Refuted: special $\lambda$ identities in every odd characteristic]
\label{C:odd_characteristic_failure}
Let $k$ have characteristic $p>2$. Exterior powers on
$WG_{\mathrm{sp}}(B_{C_p})$ do not give a special $\lambda$-ring. The failure
can be detected by the $p$-th Adams operation on two actual symmetric
representations. For every algebraically closed such $k$ there is also a
smooth projective connected curve $X/k$ for which $WG_{\mathrm{sp}}(X)$ is
not special under exterior powers.
\end{theorem}
\begin{proof}
Let $A_k(C_p)$ be the split Green ring and let $V_j$, $1\leq j\leq p$, be
the indecomposable module on which a generator has one unipotent Jordan
block of size $j$. Jordan normal form gives an additive basis
$[V_1],\ldots,[V_p]$ of this ring; put $V_0=0$. We first establish the two
Adams formulas needed, without using a modular character formula.

Write $x=[V_2]$ and $z=[V_p]$. For $1\leq j<p$ and for $j=p$, respectively,
\[
 x[V_j]=[V_{j-1}]+[V_{j+1}],\qquad xz=2z.                 \tag{C.3.1}
\]
For a direct matrix proof, write $J_j=1+t$ on $k[t]/(t^j)$. The nilpotent
part of $J_2\otimes J_j$ is conjugate to
$N=\left(\begin{smallmatrix}t&1\\0&t\end{smallmatrix}\right)$, and
\[
 N^a=\begin{pmatrix}t^a&a t^{a-1}\\0&t^a\end{pmatrix}.
\]
For $1<j<p$, its kernel has dimension two, $N^j$ has rank one, and
$N^{j+1}=0$, so its two blocks have sizes $j+1,j-1$. The case $j=1$ is
immediate. For $j=p$ its kernel has dimension two and $N^p=0$ on a
$2p$-dimensional space, so both blocks have size $p$. This proves (C.3.1).

Since $\lambda_t(V_2)=1+xt+t^2$, its Adams polynomials satisfy
$P_0=2$, $P_1=x$, and $P_a=xP_{a-1}-P_{a-2}$. Introduce
$U_0=0,U_1=1,U_{a+1}=xU_a-U_{a-1}$; then $P_a=U_{a+1}-U_{a-1}$.
Equation (C.3.1) gives $U_j=[V_j]$ for $j\leq p$ and
$U_{p+1}=2[V_p]-[V_{p-1}]$. Consequently
\[
 \psi^p(V_2)=2[V_p]-2[V_{p-1}].                         \tag{C.3.2}
\]
The regular module $V_p$ has its permutation basis. For $0<i<p$ the
$p$-cycle acts freely on its $i$-element subsets; each signed orbit in the
exterior power is a regular module, since the product of its signs is one.
The top exterior power is trivial because $p$ is odd. Hence
\[
 \lambda_t(V_p)=1+\sum_{i=1}^{p-1}\frac{\binom pi}{p}z t^i+t^p.
\]
Also $z^2=pz$. The ring $\mathbb Z[z]/(z^2-pz)$ injects into
$\mathbb Z\times\mathbb Z$ by evaluating $z$ at $0$ and $p$.
The displayed exterior polynomial becomes $1+t^p$ and $(1+t)^p$ at these
two values. Newton's formula gives $\psi^p=p$ at both, proving
\[
 \psi^p(V_p)=p[V_1].                                    \tag{C.3.3}
\]

Take the actual symmetric representations
$A=H(V_2)$ and $B=V_p$ with its invariant dot product. Their underlying
classes are $2[V_2]$ and $[V_p]$, and $[A\otimes B]=4[V_p]$.
The forgetful map from split $WG$ to the split Green ring commutes with
exterior powers and Adams operations. Adams operations are additive for any
pre-$\lambda$-structure. By (C.3.2)--(C.3.3), the image of their
multiplicativity defect is
\[
 \psi^p(A)\psi^p(B)-\psi^p(A\otimes B)
       =4p\bigl([V_p]-[V_{p-1}]-[V_1]\bigr)\ne0.        \tag{C.3.4}
\]
In a special $\lambda$-ring all Adams operations are multiplicative, so
this disproves specialness. The nonzero conclusion uses the split basis;
all these classes would instead equal their dimensions in the exact
representation group. For $p=3$, the attached program constructs the forms
and the actual exterior matrices through dimension $\binom{12}{3}=220$;
it independently obtains the coefficient vector $(-12,-12,12)$ in the
basis $(V_1,V_2,V_3)$.

To obtain schemes, let $Y\to X$ be a connected etale $C_p$-torsor of smooth
projective connected curves over an algebraically closed $k$.
Such torsors exist over every algebraically closed field of characteristic
$p$. Waterhouse, Theorem 4.1(1), p.536 \cite{C-Waterhouse1969}, gives an
elliptic curve $Y_0/\mathbb F_p$ with Frobenius trace one, hence
$|Y_0(\mathbb F_p)|=p$. Its rational points form a constant subgroup $C_p$.
The quotient isogeny $Y_0\to Y_0/C_p$ is an etale $C_p$-torsor: translation
is free, and a free action of the finite etale constant group gives a finite
etale quotient map. Base change to $k$ yields the desired $Y\to X$.
For any representations $V,W$, descent and $H^0(Y,\mathcal O_Y)=k$ give
\[
 \operatorname{Hom}_X(Y\times^{C_p}V,Y\times^{C_p}W)
       =(H^0(Y,\mathcal O_Y)\otimes_k\operatorname{Hom}_k(V,W))^{C_p}
       =\operatorname{Hom}_{C_p}(V,W).
\]
Thus the associated-bundle functor is fully faithful and tensorial.
The category of vector bundles on a proper curve is idempotent complete
and has finite-dimensional Hom spaces. It is therefore Krull--Schmidt:
Fitting's lemma makes the endomorphism ring of an indecomposable local, and
the usual exchange argument gives uniqueness of indecomposable summands.
Full faithfulness preserves distinct indecomposables and their local
endomorphism rings. Hence the induced homomorphism of split Grothendieck
groups is injective. It follows that the defect (C.3.4) stays nonzero in
the split vector-bundle ring of $X$. The invariant forms on $A,B$ descend,
and exterior powers commute with descent, proving the assertion for
$WG_{\mathrm{sp}}(X)$.
\end{proof}

\begin{corollary}[Proved: failure on the universal orthogonal representation]
\label{C:BO_positive_failure}
Let $k$ be any field of characteristic $p>2$, and let $m\geq p+4$.
The split ring $WG_{\mathrm{sp}}(B_{O_m})$ is not special under exterior
powers. Already its standard symmetric representation $E_m$ satisfies
\[
 \psi^p(E_m)^2\ne\psi^p(E_m\otimes E_m).
\]
Thus the universal orthogonal rings needed by the characteristic-zero
proof themselves fail in positive characteristic.
\end{corollary}
\begin{proof}
We retain $x=[V_2]$, $z=[V_p]$ and put $u=[V_{p-1}]$ in the split Green
ring of $C_p$. First
\[
 u^2=1+(p-2)z,\qquad (z-u)^2=1.                         \tag{C.3.7}
\]
Indeed $V_{p-1}$ is the augmentation ideal in $k[C_p]$. Tensoring its
defining exact sequence with $V_{p-1}$ gives a presentation of
$V_{p-1}$ by the free module $V_p^{\oplus(p-1)}$, with kernel
$V_{p-1}\otimes V_{p-1}$. Over the local ring $k[C_p]\simeq k[t]/t^p$,
this surjection is the direct sum of the standard surjection
$V_p\to V_{p-1}$ and a zero map from $p-2$ free summands: a vector
mapping to a generator can be completed to a basis, and elementary column
operations kill the other images. Its kernel is therefore
$V_1\oplus V_p^{\oplus(p-2)}$. This proves the first identity.
The second follows from $z^2=pz$ and $zu=(p-1)z$; tensoring the regular
module with any module is a free module of the corresponding dimension.

The three-dimensional module $V_3$ is self-dual with determinant one, so
$\lambda_t(V_3)=1+[V_3]t+[V_3]t^2+t^3$. Equation (C.3.1) says
$[V_3]=x^2-1$, also when $p=3$. The Newton polynomial for this exterior
series is
\[
 \psi^p(V_3)=P_p(x)^2-1=4(z-u)^2-1=3.
\]
For the polynomial identity in the middle, substitute $x=s+s^{-1}$:
the three formal roots are $s^2,1,s^{-2}$, whose $p$-power sum is
$(s^p+s^{-p})^2-1=P_p(x)^2-1$. This is an identity of integral
polynomials and requires no special-$\lambda$ assumption on the Green
ring. Consequently the self-product defect of $x$ vanishes:
$\psi^p(x^2)=\psi^p(1+V_3)=4=\psi^p(x)^2$.
The self-product defect of $z$ also vanishes, by $z^2=pz$ and
$\psi^p(z)=p$.

Now put $W=H(V_2)\perp V_p$, with the dot product on the latter summand.
Its underlying class is $w=2x+z$ and its dimension is $p+4$.
Expanding its self-product defect and using (C.3.4) gives
\[
 \psi^p(w)^2-\psi^p(w^2)
       =8p\bigl(z-u-1\bigr)\ne0.                     \tag{C.3.8}
\]
After extension to an algebraic closure, its symmetric form is isometric
to the split standard form. Its action therefore defines a homomorphism
$C_p\to O_{p+4}$ whose pullback of the standard representation is $W$.
Adding $m-p-4$ trivial orthogonal lines leaves this defect unchanged.
Pullback followed by forgetting forms detects the displayed nonzero
defect of $E_m$. Nonvanishing after scalar extension implies nonvanishing
over the original field $k$ as well.

The exact matrix certificate \texttt{code/C\_orthogonal\_certificate.py}
checks $p=3,5,7$. It computes all Adams operations on Jordan summands
from actual exterior-power matrices and obtains the tensor-square
decomposition from the actual matrix of $W\otimes W$. For $p=3$ the
defect is $(-24,-24,24)$ in the ordered basis $(V_1,V_2,V_3)$.
\end{proof}

\begin{lemma}[Proved: the $p$-th exterior Adams operation on every Jordan block]
\label{C:Jordan_Adams_p}
Let $k$ be a field of characteristic $p>2$, and let $V_j$ be the
indecomposable $k[C_p]$-module of dimension $j$, $1\leq j\leq p$.
In the split Green ring put $w=[V_p]-[V_{p-1}]$. Then $w^2=1$ and
\[
 \psi^p(V_j)=\begin{cases}j[V_1],&j\text{ odd},\\jw,&j\text{ even}.
 \end{cases}                                                   \tag{C.3.10}
\]
If $x=[V_2]$ and $D(a,b)=\psi^p(a)\psi^p(b)-\psi^p(ab)$, then
\[
 D(x,V_j)=0\quad(j<p),\qquad D(x,V_p)=2p(w-1).             \tag{C.3.11}
\]
\end{lemma}
\begin{proof}
The tensor calculations in the proofs of
Theorem~\ref{C:odd_characteristic_failure} and
Corollary~\ref{C:BO_positive_failure} give
$xV_j=V_{j-1}+V_{j+1}$ for $j<p$, $xV_p=2V_p$, and $w^2=1$.
These steps use actual Jordan kernels and the augmentation-ideal sequence;
they have been rechecked for their use here.
Define $U_0=0$, $U_1=1$, $U_{j+1}=xU_j-U_{j-1}$.
The first identity gives $V_j=U_j(x)$ for $j\leq p$.
Let $P_a$ be the integral polynomial determined by
$P_a(s+s^{-1})=s^a+s^{-a}$. Since $\lambda_t(V_2)=1+xt+t^2$,
Newton's recursion gives $\psi^a(x)=P_a(x)$ for every $a$.

For $a$ prime to $p$, the exterior Adams operation on a modular Green
ring is a ring endomorphism: Bryant--Johnson
\cite[Proposition 2.1]{C-BryantJohnson2011}. Applied with $1\leq a<p$,
this gives $\psi^a(V_j)=U_j(P_a(x))$.
In the torsion-free Green ring, Newton's recursions therefore determine
$\lambda^r(V_j)$ for $r<p$ to be the same integral polynomial in $x$
as the exterior character of the formal weights
$s^{j-1},s^{j-3},\ldots,s^{1-j}$. One may perform this comparison after
tensoring with $\mathbb Q$, where division by $1,\ldots,p-1$ is legitimate,
and descend because the Green ring has the free basis $V_1,\ldots,V_p$.
For $j<p$ these are every nonzero exterior power. For $j=p$ the one
remaining exterior power is the determinant $\lambda^p(V_p)=1$, also
equal to the product of those formal weights. Newton's formula at $p$
now gives
\[
 \psi^p(V_j)=U_j(P_p(x))=U_j(2w)=j w^{j-1}.
\]
Here $P_p(x)=2w$ was obtained directly from the recurrence at the boundary
$xV_p=2V_p$, and the last equality follows inductively from $w^2=1$.
This proves (C.3.10). Substitution in $xV_j=V_{j-1}+V_{j+1}$ shows
(C.3.11), for odd and even $j$ separately; at $j=p$ one instead uses
$xV_p=2V_p$ and $\psi^p(V_p)=p$.
\end{proof}

\begin{theorem}[Proved: the complete fieldwise special-$\lambda$ criterion]
\label{C:BO_field_lambda}
Let $k$ be any field of characteristic $p>2$, and let $O_m$ be the split
orthogonal group, including $O_0=1$. Then exterior powers make
$WG_{\mathrm{sp}}(B_{O_m})$ a special $\lambda$-ring if and only if
$m\leq2$. Thus the low-rank cases left by the bound $m\geq p+4$ are
all decided. The negative assertion holds over arbitrary perfect fields
and arbitrary algebraically closed fields, with no distinction in the
criterion.
\end{theorem}
\begin{proof}
For $m=0,1,2$ the representation category is semisimple. For $m=1$ this
is the eigenspace decomposition for $\mu_2$. For $m=2$ write
$O_2=\mathbb G_m\rtimes C_2$ with inversion action. Torus weights
split every representation; the zero weight splits into the two
$C_2$-characters, while each pair of nonzero opposite weights is the
sum of copies of the two-dimensional exchanged-weight representation.
This proves semisimplicity over $k$ itself. Every metabolic symmetric
representation then splits isometrically as a hyperbolic representation:
choose an equivariant complement to the Lagrangian and correct its
self-pairing by half its Gram map. Consequently the split group equals
the exact Grothendieck--Witt group. Zibrowius's main theorem, or his
Proposition 3.11 for $m=2$, applies to that group
\cite{C-Zibrowius2015}. Its field hypothesis and $2\ne0$ are exactly
our hypotheses, proving specialness.

We next treat $m\geq4$. Let $E_m$ be the standard symmetric
representation and give $B=\operatorname{Sym}^{p-1}(E_m)$ the symmetric
form obtained by restricting the tensor-power form to the image of the
symmetrizing idempotent. This idempotent exists since $(p-1)!$ is a
unit and is self-adjoint; its image is therefore nondegenerate.
Use the homomorphism $C_p\to O_m$ with representation
$H(V_2)\perp E_{m-4}^{\mathrm{triv}}$, where the second term is the
standard split quadratic space of dimension $m-4$ with trivial
$C_p$-action. It is defined over $k$: the hyperbolic form is split and
the indicated split orthogonal summand is fixed.
In the Green ring the restriction of $E_m$ is $2x+(m-4)$.
Write $b_j\geq0$ for the Jordan multiplicities of $B$ on this subgroup.
We have $b_p>0$: the direct-sum formula for symmetric powers exhibits
$\operatorname{Sym}^{p-1}(V_2)$ as a summand of the restriction, and
this module is $V_p$. Indeed on polynomials of degree at most $p-1$
the generator acts by $t\mapsto t+1$; its successive differences have
nonzero leading coefficients $1,2,\ldots,p-1$, giving one Jordan block.
The bilinearity of the defect $D$ follows from additivity of Adams
operations. Lemma~\ref{C:Jordan_Adams_p} thus gives
\[
 D(E_m,B)\big|_{C_p}=4p b_p(w-1)\ne0.                  \tag{C.3.12}
\]
Nonvanishing follows from the free indecomposable basis, since
$w-1=V_p-V_{p-1}-V_1$. Restriction and forgetting forms commute with
exterior powers. Hence (C.3.12) detects a nonzero Adams
multiplicativity defect in $WG_{\mathrm{sp}}(B_{O_m})$.

For $m=3$ use $SO_3\simeq PGL_2$, valid over every field with $2\ne0$.
Conjugation on the traceless two-by-two matrices preserves the split
quadratic form $-\det$ and has kernel the scalars. It identifies the
smooth connected three-dimensional groups $PGL_2$ and $SO_3$: its
Lie map is injective, and its image is consequently an open subgroup
of the connected target. Let $U=k^2$ be the standard $SL_2$-module and set
\[
 A=H(\operatorname{Sym}^{2p-2}U),\qquad
 B=\operatorname{Sym}^{p-1}U.
\]
The central $\mu_2$ acts trivially on both, so they descend to $PGL_2$.
The invariant alternating form on $U$, tensor-powered $p-1$ times and
restricted by the symmetrizing idempotent, gives $B$ an invariant
nondegenerate symmetric form. The hyperbolic form on $A$ is invariant
by construction. Extend both to $O_3$ along the projection
$g\mapsto\det(g)g$ from $O_3$ onto $SO_3$.

Restrict to the upper-unitriangular constant subgroup $C_p$ of $SL_2$.
On degree-at-most-$2p-2$ polynomials, write $v=t^p-t$. The basis
\[
 1,t,\ldots,t^{p-1};\qquad v,vt,\ldots,vt^{p-2}
\]
is triangular by polynomial degree. Each block is stable under
$t\mapsto t+1$, since $v$ is invariant, and has a single Jordan block
by the same difference calculation. Hence
$\operatorname{Sym}^{2p-2}U|_{C_p}=V_p\oplus V_{p-1}$ and
$B|_{C_p}=V_p$. All Jordan modules are self-dual, so
$A|_{C_p}=2V_p+2V_{p-1}$. As $V_p\otimes V_j=jV_p$,
Lemma~\ref{C:Jordan_Adams_p} gives
\[
 D(A,B)\big|_{C_p}=2p(p-1)(w-1)\ne0.                 \tag{C.3.13}
\]
This proves failure also for $m=3$.

The program \texttt{code/C\_finish\_lowrank\_lambda.py} constructs the
$O_3$ and $O_4$ witnesses for $p=3,5,7$, including their symmetric
forms, actual symmetric-power matrices and tensor-product Jordan
blocks. It independently computes the Adams operations from actual
exterior matrices on every Jordan block. The respective defect
coefficients in front of $w-1$ are $12,40,84$ for $O_3$ and
$36,100,196$ for $O_4$; the largest tensor matrix has dimension $336$.
These computations verify six finite witnesses and every matrix identity
used in their certificates; the preceding argument proves all primes
and all ranks.
\end{proof}

\begin{proposition}[Proved: the rings in ranks zero, one and two over any affine base]
\label{C:BO_lowrank_base}
Let $k$ be any commutative ring with $2\in k^\times$, and put
$W=GW(k)$, the group completion of nondegenerate symmetric projective
$k$-modules under orthogonal sum. Then
\[
 WG_{\mathrm{sp}}(B_{O_0})=W,\qquad
 WG_{\mathrm{sp}}(B_{O_1})=W[\delta]/(\delta^2-1),
\]
where $\delta$ is the sign representation with form $\langle1\rangle$.
For $O_2$, put $d=(\det,\langle-1\rangle)$ and let $u_a$, $a\geq1$,
be the exchanged-weight representation of weights $a,-a$ with Gram
matrix $\left(\begin{smallmatrix}0&1\\1&0\end{smallmatrix}\right)$.
Then the answer is the free $W$-module
\[
 W\,1\ \oplus W\,d\ \oplus\bigoplus_{a\geq1}W\,u_a,
\]
with multiplication
\[
 d^2=1,\quad du_a=u_a,\quad
 u_a u_b=u_{a+b}+u_{|a-b|},\qquad
 u_0=\langle2\rangle(1+d).                             \tag{C.3.14}
\]
For the displayed basis forms,
$\lambda_t(d)=1+dt$ and $\lambda_t(u_a)=1+u_at+dt^2$.
For a coefficient form $(P,b)$, the exterior operations are obtained by
forming the actual finite projective representation $P\otimes u_a$
with its tensor form, taking its exterior power, and decomposing by the
same weight orbits and zero-weight eigenspaces. This coefficient rule
and the orthogonal-sum formula specify every exterior operation,
without any field, noetherian, reducedness or freeness assumption on $k$.
\end{proposition}
\begin{proof}
The standard equivalence between $\mathbb G_m$-comodules and
$\mathbb Z$-graded modules is immediate from the basis $t^a$ of
$k[t,t^{-1}]$: coassociativity makes the coefficient maps orthogonal
idempotents. A finitely generated projective comodule has finitely many
nonzero weight modules, each a projective direct summand. The reflection
in $O_2=\mathbb G_m\rtimes C_2$ exchanges the $a$ and $-a$ modules.
An invariant form pairs only opposite weights, since the Laurent
monomials remain a $k$-basis. On an opposite-weight pair identify the
negative weight with the positive weight by reflection. The pairing
then becomes a symmetric nondegenerate form $b$ on that projective
module; the pair is precisely $(P,b)\otimes u_a$.
On weight zero the idempotents $(1\pm r)/2$ split the reflection;
their images are orthogonal and give the two coefficient forms.
These constructions are inverse equivalences of the corresponding
isometry groupoids under orthogonal sum. Thus group completion gives
the asserted direct sums over $W$. The same argument for $\mu_2$
gives the formula in rank one, and rank zero is the definition.

The matrices in Zibrowius, Lemmas 3.9--3.10
\cite{C-Zibrowius2015}, use only $0,1,-1$; the zero-weight change of
basis additionally has determinant $-2$. They are therefore isometries
over this arbitrary $k$, not just over a field, and yield (C.3.14).
In particular the reflection eigenvectors in the zero-weight
hyperbolic plane have norms $2,-2$, explaining the factor
$\langle2\rangle$. Direct determinants give the stated exterior
polynomials. Taking actual tensors and exterior powers commutes with
isometries and base change, while the weight/eigenspace decomposition
just proved is canonical up to isometry. This proves the coefficient
rule and its completeness.
\end{proof}

\begin{theorem}[Proved: the special-$\lambda$ criterion over every affine base]
\label{C:BO_affine_lambda}
Let $k$ be a commutative ring with $2\in k^\times$, and let $m\geq0$.
Then
\[
 WG_{\mathrm{sp}}(B_{O_m,k})\text{ is special under exterior powers}
 \quad\Longleftrightarrow\quad
       m\leq2\ \text{or}\ k\text{ is a }\mathbb Q\text{-algebra}.       \tag{C.3.15}
\]
There is no noetherian, reducedness, characteristic-purity or
finite-dimensionality hypothesis on $k$.
\end{theorem}
\begin{proof}
We first explain the exact-group input over arbitrary $k$.
Fasel--Haution, Proposition 4.2.2 and the proof of Corollary 4.2.3
\cite{C-FaselHaution2025}, establish the universal special-$\lambda$
identities in exact Grothendieck--Witt groups of $O_r$ and
$O_r\times O_s$ over $\mathbb Z[1/2]$. Their identities transfer to
exact groups of equivariant projective $k$-modules as follows.
For a symmetric $G$-representation $(P,q)$ of rank $r$, its orthogonal
frame torsor $T\to\operatorname{Spec}k$ has a left $G$-action commuting
with its right $O_r$-action. Associate to an $O_r$-representation $V$
over $\mathbb Z[1/2]$ the equivariant projective module
$T\times^{O_r}(V\otimes k)$. This is an exact tensor functor preserving
forms, isometries, Lagrangians, and exterior powers. Exactness follows
on the faithfully flat trivializing cover, where it is scalar extension
of a sequence of finite projective $\mathbb Z[1/2]$-modules; the latter
sequence is split on underlying modules. A pair of frame torsors gives
the identical functor for $O_r\times O_s$.
Thus the universal iterated-exterior and product identities hold in
the exact equivariant group. Their extension from actual objects to
virtual classes uses the pre-$\lambda$ addition formula. Finitely many
open-and-closed rank strata suffice for every pair of projectives, so
constant rank need not be assumed.

For $m\leq2$, every admissible sequence of $O_m$-representations
splits equivariantly. In ranks zero and one this follows from
projectivity and the two sign idempotents; in rank two use the weight
and reflection decomposition in Proposition~\ref{C:BO_lowrank_base}.
On each multiplicity module the quotient is projective, so its sequence
splits over $k$, and these splittings reconstruct an equivariant one.
Any equivariant Lagrangian consequently has an equivariant complement,
and the usual correction by half its Gram map makes that complement
isotropic. Metabolic relations are therefore already actual isometries,
and hyperbolic exact-sequence relations are already orthogonal-sum
isometries. The split and exact groups coincide. The preceding
universal identities prove specialness.

If $k$ is a $\mathbb Q$-algebra, the same conclusion holds for every
$m$. The coordinate coalgebra $\mathbb Q[O_m]$ is a direct sum of
matrix coalgebras, since its finite-dimensional representation category
is split semisimple. After tensoring with $k$, every finite projective
comodule has a finite isotypic decomposition, whose multiplicity modules
are projective direct summands. A sequence therefore splits on each
multiplicity module and hence equivariantly. The preceding comparison
again identifies the split and exact groups and proves specialness.

Conversely assume $m\geq3$ and $k$ is not a $\mathbb Q$-algebra.
Some prime $p$ is not a unit of $k$; since $2$ is a unit, $p>2$.
Choose a maximal ideal containing $p$ and let $F$ be its residue field.
We now lift the negative witnesses to $\mathbb Z[1/2]$, so that their
nonvanishing over $F$ detects nonvanishing over $k$.
For $m\geq4$ take
\[
 E_m,\qquad H(\operatorname{Sym}^{p-1}E_m).
\]
Both are nondegenerate symmetric representations over $\mathbb Z[1/2]$;
no factorial is inverted to form the hyperbolic object. Over $F$ the
symmetric power is self-dual by the form constructed in
Theorem~\ref{C:BO_field_lambda}, so the Green-ring defect is twice
(C.3.12) and is nonzero. For $m=3$ the even-degree representations
$\operatorname{Sym}^{2p-2}(\mathbb Z[1/2]^2)$ and
$\operatorname{Sym}^{p-1}(\mathbb Z[1/2]^2)$ descend to $PGL_2$
over $\mathbb Z[1/2]$. Extend them to $O_3$ as before and take the
hyperbolic forms on both. Over $F$ their defect is twice (C.3.13),
again nonzero. Scalar extension $k\to F$ is a homomorphism of split
pre-$\lambda$-rings, since it preserves isometries, sums, tensors and
exterior powers. A zero defect over $k$ would have zero image over
$F$, contradiction. This proves necessity and the equivalence.
\end{proof}

\begin{proposition}[Proved: explicit rings over every rational affine base]
\label{C:BO_Qalgebra}
Let $k$ be any commutative $\mathbb Q$-algebra. For split $O_m$ choose
the rational simple representations $V_\lambda$ and rational invariant
symmetric forms $q_\lambda$ indexed as in
Proposition~\ref{C:orthogonal_representation_ring}. Then
\[
 WG_{\mathrm{sp}}(B_{O_m,k})
   =\bigoplus_{\lambda\in\operatorname{Irr}(O_{m,\mathbb Q})}
          GW(k)[V_\lambda\otimes k,q_\lambda\otimes k].          \tag{C.3.16}
\]
The tensor coefficients are the rational symmetric multiplicity forms
of Proposition~\ref{C:orthogonal_representation_ring}, extended to $k$.
For exterior operations one forms the actual exterior power of a finite
sum $(V_\lambda\otimes P_\lambda,q_\lambda\otimes b_\lambda)$ and
extracts its projective isotypic multiplicity forms. This gives the
whole ring and all its exterior operations over arbitrary nonfield,
nonreduced and nonnoetherian rational affine bases. It is special by
Theorem~\ref{C:BO_affine_lambda}.
\end{proposition}
\begin{proof}
Use the matrix-coalgebra decomposition from the preceding proof.
Concretely, a finite projective representation has coaction through a
finite direct sum of these coalgebras; dualizing that finite sum gives
a product of matrix algebras over $k$. The matrix idempotents exhibit
each multiplicity module $P_\lambda$ as a finite projective direct
summand and give an actual decomposition
$M=\bigoplus_\lambda V_\lambda\otimes P_\lambda$.
The simple rational $O_m$-modules have scalar endomorphisms and are
symmetric and self-dual, by the opened Zibrowius Proposition 3.17.
Consequently pairings between different components vanish, and the
pairing on one component is uniquely $q_\lambda\otimes b_\lambda$.
Nondegeneracy is equivalent to that of $b_\lambda$. This proves the
isometry classification and its group completion (C.3.16).
Tensor decompositions over $\mathbb Q$ extend to $k$ and tensor the
multiplicity forms with $b_\lambda\otimes b_\mu$. Exterior powers
remain finite projective representations, so the same canonical
isotypic extraction applies. This proves the asserted operations and
that no additional additive relations or objects are omitted.
\end{proof}

\begin{proposition}[Proved reduction: the remaining rank-three module problem contains a wild category]
\label{C:BO3_wild_gate}
Let $k$ be algebraically closed of characteristic $p>2$, and put
$d=p^2+1$ if $p\geq5$, and $d=12$ if $p=3$. There is a fully faithful
exact functor
\[
 F:\operatorname{mod}S_k(2,d)\longrightarrow
             \operatorname{Rep}^{\mathrm{fd}}_k(O_3),
\]
where $S_k(2,d)$ is a Schur algebra of wild representation type.
The functor preserves indecomposables. With
$M^\dagger=M^*\otimes\det^d$, viewed as a degree-$d$ polynomial
$GL_2$-module, for indecomposable $M,N$ one has
\[
 [H(FM)]=[H(FN)]\text{ in }WG_{\mathrm{sp}}(BO_3)
 \quad\Longleftrightarrow\quad N\simeq M\text{ or }N\simeq M^\dagger.
                                                               \tag{C.3.17}
\]
Thus the additive split group retains these wild modules up to the
specified duality. This is a rigorous reduction of the outstanding
classification problem, not an enumeration of the indecomposables
and not a completed answer to the full modular-ring question.
\end{proposition}
\begin{proof}
The coalgebra of homogeneous degree-$d$ polynomials in the four matrix
coordinates has dual $S_k(2,d)$. Finite-dimensional modules over its
dual are exactly its finite-dimensional comodules, hence homogeneous
polynomial $GL_2$-representations of degree $d$. This equivalence,
including the indicated scalar weight, is stated in
Doty--Erdmann--Martin--Nakano, Section 1.3, printed p.140
\cite{C-DotyErdmannMartinNakano1999}.
Their Proposition in Section 3.4, printed pp.151--152, says that
$S_k(2,d)$ is wild for $p\geq5,d\geq p^2$, and for $p=3,d\geq12$.
The precise cases are confirmed by
Hemmer--Kujawa--Nakano, Theorem 1.2.1(b)--(d)
\cite{C-HemmerKujawaNakano2006}. The original authors' correction for
$S(2,11)$ in characteristic two does not affect either chosen case.

On every homogeneous degree-$d$ module scalar matrices act by $t^d$.
Therefore a linear map between two such modules commuting with $SL_2$
also commutes with scalar matrices. Over algebraically closed $k$ every
$g\in GL_2(k)$ is a scalar times an element of $SL_2(k)$, by choosing
a square root of $\det(g)$. Such a map is consequently $GL_2$-equivariant;
the equality of rational morphisms follows from density of $k$-points
in the smooth reduced group. Restriction to $SL_2$ is fully faithful.
Since $d$ is even, the central $\mu_2$ acts trivially, so these modules
descend to $PGL_2\simeq SO_3$. Extend to $O_3$ along
$g\mapsto\det(g)g$. All these functors are exact on vector spaces and
fully faithful, proving the assertion about $F$. Its equality of
endomorphism rings proves preservation of indecomposables.

For clarity $M^\dagger$ is again polynomial of degree $d$:
the inverse matrix formula has denominator $\det(g)$, so the degree-$d$
matrix coefficients of the dual become polynomial after multiplication
by $\det(g)^d$, with total degree $d$. On $SL_2$ that determinant
twist is trivial, giving $F(M^\dagger)\simeq F(M)^*$.
If the two classes in (C.3.17) agree, forgetting forms gives
$[FM]+[(FM)^*]=[FN]+[(FN)^*]$ in the split Green group.
Krull--Schmidt and the free indecomposable basis give
$FN\simeq FM$ or $FN\simeq(FM)^*$. Full faithfulness gives the two
asserted possibilities. Conversely an isomorphism of modules induces
an isometry of their hyperbolic forms, and $H(V)\simeq H(V^*)$ by
exchanging its two summands. This proves the reverse implication and
specifies exactly the ambiguity retained by hyperbolic passage.
\end{proof}


\begin{proposition}[Known structure theorem applied to the split ring over any field]
\label{C:field_WG_structure}
Let $k$ be a field with $2\in k^\times$ and $G$ an affine group scheme.
For every self-dual indecomposable finite-dimensional representation $U$,
choose an invariant nondegenerate $\varepsilon_U$-symmetric form $q_U$,
where $\varepsilon_U\in\{1,-1\}$. Such a choice always exists.
Put
\[
 E_U=\operatorname{End}_G(U),\qquad
 D_U=E_U/\operatorname{rad}(E_U),
\]
and let $\sigma_U$ be the involution on the division algebra $D_U$
induced by the $q_U$-adjoint involution. Then, as an abelian group,
\[
 WG_{\mathrm{sp}}(B_G)\simeq
 \bigoplus_{U\simeq U^*}GW^{\varepsilon_U}(D_U,\sigma_U)
 \ \oplus\!
 \bigoplus_{\{U,U^*\}:U\not\simeq U^*}\mathbb Z[H(U)].       \tag{C.3.9}
\]
Here $GW^{\varepsilon}(D,\sigma)$ means the group completion under
orthogonal sum of nondegenerate $\varepsilon$-hermitian forms on finite
right $D$-vector spaces; it does not mean a quotient by hyperbolic forms.

The product and exterior operations are determined by the following rule.
Lift a residue multiplicity form to the corresponding isotypic
representation $U^{\oplus r}$, take the tensor product or exterior power
there, decompose the resulting actual representation by indecomposable
duality orbits, and reduce the induced multiplicity forms modulo the
endomorphism radicals. The result is independent of the lifts. This
applies in particular to $G=O_m$ over every such field.
The radical quotient of the representation category need not be a tensor
quotient: tensor products in this rule are taken before reduction.
\end{proposition}
\begin{proof}
This is the specialization of Quebbemann--Scharlau--Schulte
\cite[Propositions 2.4--2.5 and Theorems 3.2--3.3, pp.271--273]{C-QSS1979}
to the additive category of finite-dimensional $G$-representations with
contragredient duality. All idempotents split; each object has a
Krull--Schmidt decomposition; and its endomorphism algebra is finite
dimensional over $k$, hence complete for its nilpotent radical. These
verify all three hypotheses in their \S3. Since two is invertible,
quadratic and symmetric forms are identified by symmetrization, as in
their \S1, p.266.

Proposition 2.5 supplies $q_U$. Transfer in Proposition 2.4 identifies
forms on sums of $U$ with $\varepsilon_U$-hermitian multiplicity forms
over $E_U$. Reduction modulo its radical preserves isometry classes
by Theorem 2.2 and Corollary 2.3, pp.269--270. Theorem 3.3 separates
the duality orbits uniquely and makes each non-self-dual orbit
hyperbolic. Group completion yields (C.3.9), with a free integer
coefficient recording the number of copies of $U$ in that last case.
The stated operations use actual tensors and exterior powers, both of
which preserve isometries. The same reduction theorem shows independence
of the chosen lifts, proving the ring description with its operations.
These are the split $KQ_0$ groups defined on p.267 of the source; its later
exact $GQ_0$ groups are not being substituted.
\end{proof}

\begin{proposition}[Known structure theorem specialized to an algebraically closed field]
\label{C:algebraically_closed_WG}
Let $k$ be algebraically closed with $\operatorname{char}k\ne2$, and let
$G$ be an affine group scheme over $k$. Write $A_k(G)$ for the split
Green ring of finite-dimensional representations. Forgetting the form
embeds $WG_{\mathrm{sp}}(B_G)$ into $A_k(G)$, compatibly with products
and exterior powers. Its additive basis is as follows:
\begin{enumerate}
\item $[U]$ for every self-dual indecomposable $U$ admitting a nondegenerate
invariant symmetric form;
\item $2[U]$ for every self-dual indecomposable $U$ admitting a nondegenerate
invariant alternating form;
\item $[U]+[U^*]$ for every unordered pair of nonisomorphic dual
indecomposables.
\end{enumerate}
Every self-dual indecomposable has exactly one of the first two types.
This determines the ring, including its exterior operations, as a specified
subring of the Green ring. In particular it applies to $G=O_m$ in every
odd characteristic; it does not assert that the indecomposable
$O_m$-modules have a semisimple highest-weight classification.
\end{proposition}
\begin{proof}
Apply Proposition~\ref{C:field_WG_structure}. Every finite-dimensional
division algebra over an algebraically closed field is $k$, since each
element has a split minimal polynomial and hence is scalar. The residual
adjoint involutions are therefore the identity. Proposition 2.5(2) of
\cite{C-QSS1979} shows that the symmetric and alternating types cannot
both occur on one indecomposable. Symmetric multiplicity forms over $k$
are classified by their dimension; alternating multiplicity forms are
classified by their even dimension. The three summands in (C.3.9)
therefore have precisely the underlying-module generators stated above.
Their supports in the free abelian indecomposable basis are disjoint,
which proves injectivity. Isometry being determined by the underlying
module is also exactly Quebbemann--R.~Scharlau--W.~Scharlau--Schulte,
\cite[Corollary 3.5, p.97]{C-QSSS1976}, or
\cite[Application 3.4(3), p.274]{C-QSS1979}.
Forgetting forms preserves tensors and exterior powers, giving the
remaining assertions.
\end{proof}

\begin{proposition}[Proved: the hyperbolic Stiefel--Whitney image]
\label{C:hyperbolic_SW}
Let $S$ be a scheme with $2$ invertible, and use the etale Stiefel--Whitney
classes with rank-one normalization. Put $a=\partial(-1)\in
H^1_{\mathrm{et}}(S,\mathbb Z/2)$. For a vector bundle $E$ of rank $r$,
\[
 w(H(E))=\sum_{j=0}^r(1+a)^{r-j}\overline c_j(E).
                                                               \tag{C.3.5}
\]
Here $\overline c_j(E)\in H^{2j}_{\mathrm{et}}(S,\mathbb Z/2)$ is the
mod-two Chern class, using $\mu_2\simeq\mathbb Z/2$ to identify twists.
Consequently the image of the hyperbolic map in the target
$H^0(S,\mathbb Z)\times(1+\prod_{i>0}H^i_{\mathrm{et}}(S,\mathbb Z/2))$
is precisely the set of pairs
\[
 \left(2(r-s),\,
 \frac{\sum_{j=0}^r(1+a)^{r-j}\overline c_j(E)}
      {\sum_{j=0}^s(1+a)^{s-j}\overline c_j(F)}\right),          \tag{C.3.6}
\]
for vector bundles $E,F$ of ranks $r,s$, componentwise for nonconstant rank.
The formula factors through exact $K_0(S)$, even when its source is split
$K_s(S)$. In particular the image is not in general trivial in the
Stiefel--Whitney coordinate.
\end{proposition}
\begin{proof}
First let $L$ be a line bundle. The scheme $P'$ of nonisotropic lines in
$H(L)$ is the complement of the zero section in $L^{\otimes2}$ (replacing
this line bundle by its dual gives the same torsor). Its tautological
quadratic line has underlying bundle $L_{P'}$ and a class
$y\in H^1(P',\mu_2)$. Its orthogonal complement has class $y+a$.
Thus Whitney's formula gives
\[
 w(H(L))|_{P'}=(1+y)(1+y+a)=1+a+y^2+ay.
\]
The connecting map for $1\to\mu_2\to\mu_4\to\mu_2\to1$ sends
$y$ to $y^2+ay$. One can see the correction term directly: on a common
etale cover choose a square root $i$ of $-1$; writing a $\mu_2$ cocycle
as $(-1)^{y_{ij}}$, lift it to $i^{y_{ij}}$. The coboundary has exponent
$y_{ij}y_{jk}+a_{ij}y_{jk}$ modulo two. The same connecting map is the
Kummer Chern class of the underlying line bundle. Therefore
$y^2+ay=\overline c_1(L_{P'})$.

Pullback $H^*(S,\mathbb Z/2)\to H^*(P',\mathbb Z/2)$ is injective: this
is the rank-two instance of the nonisotropic-projective-bundle theorem
used in Laborde, \S III, p.50 \cite{C-Laborde1976}. Hence
$w(H(L))=1+a+\overline c_1(L)$ on $S$.
For general $E$, the usual splitting principle pulls it to a direct sum of
line bundles by an iterated projective bundle followed by the affine
bundle of splittings. Pullback in mod-two etale cohomology is injective
by the projective-bundle formula and affine-bundle invariance. Applying the
rank-one formula and Whitney's formula there gives
$\prod_i(1+a+\xi_i)$, which expands to (C.3.5) and descends by injectivity.

The total Stiefel--Whitney class is multiplicative on orthogonal sums, so
it extends to the group completion by ratios; each series with constant
term one has a unique inverse in the degree completion. This proves both
inclusions in (C.3.6). Chern classes and rank satisfy the exact-sequence
Whitney relation, proving the factorization through exact $K_0$.
For example $w_1(H(E))=ra$ and
$w_2(H(E))=\overline c_1(E)+\binom r2a^2$.
The supplied study, reading 133\_002--018\_030--054, lines 80--124,
already derives the rank-one formula and credits Michel Hua's project for
linking it to the cancellation of $w\circ\alpha=1$ on scan 9. We credit
that derivation and link; the argument above rechecks the etale connecting
map and Laborde's precise splitting input.
\end{proof}

\subsection{C.4: central extensions and the formal splitting assertions}

\begin{theorem}[Proved: all central-extension twists satisfying both tests]
\label{C:central_twists}
Let $G_0$ be an ordinary group, with specified normal subgroup $N$ and
specified subgroup $N'$, such that
\[
 N'\subseteq Z(G_0),\qquad G_0'\subseteq N.
\]
Set $C=N\cap N'$, $Q=G_0/(NN')$, and use the commutator convention
$[g,h]=ghg^{-1}h^{-1}$. Both $C$ and $Q$ are abelian. Let
$E$ be a central extension $1\to C\to E\to Q\to1$, and let $G_E$ be the
corresponding central twist of $G_0$. Define
\[
 P=C_{G_0}(N)\cap\{g\in G_0:gC\in Z(G_0/C)\},\quad
 \zeta=P/N',\quad \eta=N/[G_0,N].
\]
The groups $\zeta$ and $\eta$ are abelian. Write $\alpha:\zeta\to Q$ and $\beta:C\to\eta$
for the natural maps. There are well-defined homomorphisms
\[
 \begin{aligned}
 \psi_0&:\zeta\longrightarrow\operatorname{Hom}(Q,C),
       &\psi_0(gN')(hNN')&=[g,h],\\
 \lambda_0&:\Lambda^2_{\mathbb Z}Q\longrightarrow\eta,
       &\lambda_0(\bar g\wedge\bar h)&=[g,h]\pmod{[G_0,N]}.
 \end{aligned}
\]
Let $b_E:\Lambda^2Q\to C$ be the commutator pairing of the central
extension $E$. Then
\[
 \boxed{\quad
 \begin{aligned}
 N'=Z(G_E)&\iff
   \psi_0+\bigl[z\mapsto b_E(\alpha(z)\wedge-)\bigr]
          \text{ is injective},\\
 N=G_E'&\iff \lambda_0+\beta\circ b_E
          \text{ is surjective}.
 \end{aligned}\quad}                                      \tag{C.4.1}
\]
Every alternating biadditive pairing $b:\Lambda^2Q\to C$ occurs as $b_E$.
For each such $b$, its central extensions form a torsor under
$\operatorname{Ext}^1_{\mathbb Z}(Q,C)$. Consequently (C.4.1), followed by
these torsors over precisely the admissible $b$, determines all central
extensions in the stated twisting family that satisfy both requirements.
Nonabelian $\pi_1$ cannot occur under $N'\subseteq Z(G)$; nonabelian
$\pi_0$ cannot occur under $G'\subseteq N$.
\end{theorem}
\begin{proof}
Choose a normalized cocycle $e:Q\times Q\to C$ representing $E$. On the
underlying set of $G_0$ put
\[
 g*h=e(\bar g,\bar h)gh.
\]
Centrality of $C$ and the cocycle identity give associativity; normalization
gives the identity, and inverses follow directly. Cohomologous cocycles give
isomorphic marked groups by multiplying each lift by the corresponding
one-cochain. On $NN'$ the product and the conjugation action of every element
remain unchanged, because its image in $Q$ is zero. The group $G_E/C$ is
also unchanged. These observations show that $P,\zeta,\eta,\alpha,\beta$
are independent of the twist.

For $g\in P$, all $[g,h]$ lie in $C$, and $g$ commutes with $N$ and $N'$.
The central commutator identities therefore show that $h\mapsto[g,h]$
is a homomorphism factoring through $Q$, and that it is independent of the
lift of $gN'$. They also show additivity in $gN'$. Since $[P,P]\subset C$
and $C\subset N'$, the group $\zeta$ is abelian. Any central element of
$G_E$ lies in $P$, and its class lies in the kernel of the resulting map
$\psi_E$. Conversely a class in that kernel has a central lift. Thus
$Z(G_E)/N'=\ker\psi_E$.

The quotient $G_0/[G_0,N]$ has $\eta$ as a central subgroup and has abelian
quotient $G_0/N$. Its commutator therefore defines an alternating biadditive
map on $G_0/N$. As $N'$ is central, this factors through $Q$ in both
arguments, giving $\lambda_0$. Its image is exactly
$G_0'/[G_0,N]$, since commutators generate $G_0'$. Therefore
$N/G_0'=\operatorname{coker}\lambda_0$, and the same argument applies after
twisting. Directly from the twisted product,
\[
 [g,h]_E=[g,h]_0\,e(\bar g,\bar h)e(\bar h,\bar g)^{-1}
        =[g,h]_0\,b_E(\bar g\wedge\bar h).
\]
This fixes the sign left open on scan 49 and gives both formulas in
(C.4.1), as well as their necessity and sufficiency.

For completeness, surjectivity onto alternating pairings does not require
$Q$ or $C$ to be finite. Choose a free abelian presentation $F\twoheadrightarrow
Q$ with kernel $R$, order a basis of $F$, and lift $b$ to $F$. An upper
triangular bilinear cocycle on that basis gives a central extension
$E_F$ of $F$ by $C$ having this commutator pairing. The preimage of $R$
is central and abelian, since the pairing vanishes on $R\times F$.
The group $R$, being a subgroup of a free abelian group, is free abelian.
Its abelian extension by $C$ thus splits; quotienting $E_F$ by a chosen
lift of $R$ gives the desired extension of $Q$. Finally a central extension
has zero commutator pairing exactly when its middle group is abelian.
The Baer difference of two extensions with the same pairing is consequently
an abelian extension, and abelian extensions are classified by
$\operatorname{Ext}^1_{\mathbb Z}(Q,C)$. This action is free and transitive.

The corrected definition of $\zeta$ agrees with the inserted
$Z(G/\pi_1)/(N'/\pi_1)$ in scan 47, intersected with $\ker\theta$; using
$Z(G/N')$ alone would allow commutators outside $\pi_1$. The formulas
rederive the supplied study's ordinary-group analysis of scans 47--54;
they also make explicit the alternating-pairing fibers of the answer.
\end{proof}

\begin{proposition}[Refuted: crossed-module data alone determine Hall isoclinism]
\label{C:isoclinism_counterexample}
The pair of crossed modules coming from the derived subgroup and the center,
even with their identified kernels and cokernels, does not determine a
Hall-isoclinism class. Thus L9 is false if it identifies the manuscript's
crossed-module data with Hall's equivalence relation. The precise connection
is a forgetful one: isoclinism retains the commutator pairing in addition to
the data forgotten by the crossed modules.
\end{proposition}
\begin{proof}
Hall isoclinism (\cite{C-Hall1940}, \S2, p.133, conditions (1)--(3)) consists of isomorphisms $G/Z(G)\simeq H/Z(H)$ and
$G'\simeq H'$ that carry the commutator map of $G$ to that of $H$.
Consider, in characteristic three,
\[
 G_1=\operatorname{Heis}(\mathbb F_3)\times
          \operatorname{Heis}(\mathbb F_3),\qquad
 G_2=\operatorname{Heis}(\mathbb F_9).
\]
Both have order $3^6$, center equal to derived subgroup
$W\simeq\mathbb F_3^2$, and central quotient $V\simeq\mathbb F_3^4$.
Their derived crossed modules are therefore both the zero map
$W\to V$ with trivial action; their central crossed modules are the same
zero map, and the identified kernels and cokernels coincide.

A completely explicit construction is $V\times W$ with multiplication
\[
 (v,z)(w,t)=(v+w,z+t+\tfrac12 b(v,w)).
\]
For $G_1$, take
\[
 b_1(v,w)=(v_1w_2-v_2w_1,\ v_3w_4-v_4w_3).
\]
For $G_2$, identify $V=\mathbb F_9^2$, $W=\mathbb F_9$, with
$\mathbb F_9=\mathbb F_3[a]/(a^2+1)$, and take
$b_2(v,w)=v_1w_2-v_2w_1$ over $\mathbb F_9$.
Bilinearity proves associativity, and the commutator is exactly $b_i$.
Both pairings have zero radical and span $W$, proving the center and derived
claims above. For $b_1$, among the eighty nonzero $v$, the linear map
$b_1(v,-):V\to W$ has rank one for sixteen $v$ and rank two for sixty-four.
Indeed rank one means exactly one of $(v_1,v_2),(v_3,v_4)$ is nonzero.
For $b_2$, the determinant map is surjective onto $\mathbb F_9$ for every
nonzero $v$, so all eighty ranks are two. Isoclinism preserves this rank
distribution. Hence the groups are not isoclinic.

The attached finite certificate enumerates the pairings and verifies their
bilinearity, radicals, images, and rank distributions. It also checks all
eight normalized bilinear twists of $Q=C_2^2$ by $C=C_2$: exactly the four
with nonzero alternating pairing have center and derived subgroup equal to
$C$. These small tests independently check both parts of (C.4.1).
\end{proof}

\begin{proposition}[Refuted: perfect-field formal splitting, even up to isogeny]
\label{C:formal_Frobenius_counterexample}
For every prime $p$, including $p=2$, there is an exact sequence of connected
$p$-divisible groups over the perfect field $\mathbb F_p$,
\[
 0\longrightarrow\mu_{p^\infty}\longrightarrow G
       \longrightarrow\mu_{p^\infty}\longrightarrow0,          \tag{C.4.2}
\]
whose associated formal Lie groups are smooth, but which does not split
even up to isogeny. In fact $G$ is not isogenous to
$\mu_{p^\infty}^2$ as an unmarked group. This refutes both (I), with its
``smooth kernel'' reading, and (II) of scan 36 over perfect fields in the
scope requested in C.4. It does not refute the narrower alternative in (I)
where the kernel is required to be unipotent of bounded exponent.
\end{proposition}
\begin{proof}
The absolute Galois group of $\mathbb F_p$ is $\widehat{\mathbb Z}$.
Give $T=\mathbb Z_p^2$ the continuous action which sends arithmetic
Frobenius to
\[
 U=\begin{pmatrix}1&1\\0&1\end{pmatrix}.
\]
Continuity follows from $U^a=\left(\begin{smallmatrix}1&a\\0&1\end{smallmatrix}\right)$:
the action factors through the continuous projection
$\widehat{\mathbb Z}\to\mathbb Z_p$. More concretely, on
$T/p^nT$ it factors through the cyclic quotient of order $p^n$.
Finite etale descent gives an etale finite group scheme $H_n$ with geometric
points $T/p^nT$; the maps induced by multiplication by $p$ assemble these
into an etale $p$-divisible group $H$. The submodule $\mathbb Z_pe_1$ and
the quotient generated by $e_2$ have trivial action, giving
\[
 0\to\mathbb Q_p/\mathbb Z_p\to H\to\mathbb Q_p/\mathbb Z_p\to0.
\]
This uses precisely the equivalence between finite etale group schemes and
continuous finite Galois modules in Tate \cite{C-Tate1967}, \S1.4,
pp.160--161; compatibility at the displayed finite levels supplies the
$p$-divisible group directly. Cartier duality, which preserves exact
sequences (Tate, \S1.2--1.3), gives (C.4.2).

Over the algebraic closure the dual is $\mu_{p^\infty}^2$, so $G$ is
connected. Tate's Proposition 1, \S2.2, p.162, identifies connected
$p$-divisible groups over a field of characteristic $p$ with divisible
commutative formal Lie groups. These formal groups have underlying formal
schemes $\operatorname{Spf}k[[x_1,\ldots,x_d]]$, hence are formally smooth.
Both endpoints and the middle of (C.4.2) therefore satisfy the smoothness
qualifications in scan 36.

An isogeny would become an isomorphism on rational Tate modules after
Cartier duality. But $U$ is not the identity on $T\otimes\mathbb Q_p$:
its fixed space is one-dimensional, whereas the trivial two-dimensional
representation has a two-dimensional fixed space. Thus no such isogeny
exists. In particular the extension cannot split in the isogeny category.
Equivalently there is no invariant complement to $\mathbb Q_pe_1$, since
$(U-1)(e_2+ce_1)=e_1$ for every $c\in\mathbb Q_p$.
This argument is characteristic independent and covers $p=2$.
\end{proof}

\begin{proposition}[Proved: two corrected formal splitting statements]
\label{C:formal_corrected}
For commutative extensions, the following are valid replacements in two
precise subclasses.
\begin{enumerate}
\item If $0\to U\to F\to P\to0$ is an extension of fppf sheaves of
commutative groups, $P$ is $p$-divisible, and $p^NU=0$, its pullback along
$[p^N]:P\to P$ splits. Thus it splits up to the explicitly specified
isogeny of the quotient, whenever the original extension and this pullback
belong to the formal category under consideration.
\item An extension of multiplicative-type $p$-divisible groups over a
perfect field $k$ splits up to isogeny if and only if the corresponding
extension of continuous finite-dimensional $\mathbb Q_p$-representations
of $\operatorname{Gal}(k^{\mathrm{sep}}/k)$ splits. Over an algebraically
closed field this always happens; over $\mathbb F_p$ Proposition
\ref{C:formal_Frobenius_counterexample} shows that it need not happen.
\end{enumerate}
\end{proposition}
\begin{proof}
In the abelian category of fppf abelian sheaves, the group
$\operatorname{Ext}^1(P,U)$ is additive in both variables. Pullback by
$[p^N]$ on $P$ and pushout by $[p^N]$ on $U$ both multiply the extension
class by $p^N$. The latter is zero, proving (1). A zero extension class
means an actual sheaf-theoretic section; the Yoneda embedding is fully
faithful, so it is a morphism of formal groups when those sheaves are
represented by the specified formal groups.

For (2), Cartier duality identifies multiplicative-type $p$-divisible
groups with etale $p$-divisible groups, and finite etale descent identifies
the latter with finite free continuous $\mathbb Z_p$-Galois modules, by
passing through their finite levels as in the previous proof. Inverting
isogenies is exactly tensoring the Hom groups and the modules with
$\mathbb Q_p$: a Galois-equivariant rational map has bounded denominator
and becomes an integral morphism after multiplication by a power of $p$.
Consequently a rational splitting exists on one side precisely when it
exists on the other. For algebraically closed $k$ the Galois group is
trivial and vector-space sequences split.
\end{proof}
